\documentclass[11pt]{article}
\usepackage[margin=0.7in]{geometry}
\usepackage{amsmath,amssymb,amsthm}
\usepackage{booktabs}
\usepackage{enumitem}
\usepackage[colorlinks=true,linkcolor=blue,citecolor=blue,urlcolor=blue]{hyperref}

\numberwithin{equation}{section}
\newtheorem{theorem}{Theorem}[section]
\newtheorem{proposition}[theorem]{Proposition}
\newtheorem{lemma}[theorem]{Lemma}
\newtheorem{corollary}[theorem]{Corollary}
\newtheorem{conjecture}[theorem]{Conjecture}
\newtheorem*{thmA}{Theorem A}
\newtheorem*{thmB}{Theorem B}
\newtheorem*{cor*}{Corollary}
\theoremstyle{remark}
\newtheorem{remark}[theorem]{Remark}

\newcommand{\PP}{\mathbb{P}}
\newcommand{\CC}{\mathbb{C}}
\newcommand{\QQ}{\mathbb{Q}}
\newcommand{\NL}{\mathrm{NL}}
\newcommand{\VZ}{V_Z}
\newcommand{\cD}{\mathcal{D}}
\newcommand{\cM}{\mathcal{M}}
\newcommand{\cI}{\mathcal{I}}
\newcommand{\cH}{\mathcal{H}}
\newcommand{\Res}{\operatorname{Res}}
\newcommand{\Hess}{\operatorname{Hess}}
\newcommand{\codim}{\operatorname{codim}}
\newcommand{\adj}{\operatorname{adj}}
\newcommand{\Ob}{\mathrm{Ob}}
\newcommand{\gt}{\tilde g}
\newcommand{\rg}{\raggedright}
\newcommand{\run}[1]{\medskip\noindent\textbf{#1}}
\setlist{nosep,leftmargin=2em}
\emergencystretch=1.5em
\linespread{0.98}
\setlength{\textfloatsep}{12pt plus 2pt minus 3pt}

\title{\vspace{-3em}Noether--Lefschetz loci of sums of two disjoint\\ complete-intersection curves\\[4pt]
\large with an application to Movasati's proposed counterexample\\ to Harris's conjecture for octic surfaces}
\author{[author names to be added]}
\date{21 November 2026 (version 4)}

\begin{document}
\maketitle

\begin{center}
\fbox{\parbox{0.93\textwidth}{\small\textbf{Notice.} This manuscript was produced by an AI system (Claude)
during one working session, at the direction of a human user; no human mathematician has verified anything in
it. Each theorem of version~1 was checked by one independent AI referee agent (``no fatal error''), and
version~1 as a whole by a third (``referee~8'': no fatal error; sixteen required corrections, applied or superseded here).
The parts that are new in version~2---Theorem~A(6) with \S\ref{sec:F}; Lemma~\ref{lem:L},
Corollary~\ref{cor:L} and the present form of Theorem~B(4)---were each read by exactly one AI referee in a
follow-up, as working notes; Lemma~\ref{lem:branch}, on which the last clauses of Theorems~A(4) and~B(3) and
the word ``branch'' in the abstract and the Corollary rest, was written by referee~8 and read by one later AI
referee (referee~9) only; version~2 was read cold by a further AI referee (``referee~9'': no fatal error;
thirteen required corrections, applied here); version~3 itself has not been re-read, and version~4 differs
from it only in the last item of \S\ref{sec:open}. The Lean~4 kernel certifies the ring
identities behind the family at $r=-1$ (Lemma~\ref{lem:lean}) and one finite computation at one explicit sextic
(\S\ref{sec:comp}), not the theorems. Computations are exact over $\QQ$ or modulo primes, as marked in
Tables~\ref{tab:status} and~\ref{tab:types}. Several references were read only through machine extraction, or
not opened; this is marked at each citation and in the bibliography.}}
\end{center}

\begin{abstract}
Let $Z_1,Z_2\subset\PP^3$ be disjoint complete-intersection curves of types $(a_1,a_2)$, $(e_1,e_2)$ on a smooth
surface $X=\{f^{\mathsf T}Mh=0\}$ of degree $d=a_1+a_2+e_1+e_2$, and $\VZ$ the locus of such surfaces. For every
rational $r\neq0$ the Noether--Lefschetz locus $\NL(r)$ of $[Z_1]+r[Z_2]$ has exactly one excess tangent
direction over $\VZ$; its second-order term is $\frac{r+1}{r}\,A$ with $A$ independent of $r$; and $\NL(-1)$ is
smooth, a branch of the closure of the family of surfaces containing a complete intersection of type
$(a_1+e_1,a_2+e_1)$. Where $A\neq0$, $\NL(r)$ is $\VZ$ with a double structure for every $r\notin\{0,-1\}$. Under
a hypothesis on one of the curves (proved for lines, smooth conics, smooth plane curves in a range of degrees
and smooth $(2,2)$ curves) we prove
$A(\sum f_ih_jm_{ij})=\frac1{2\kappa_1}\Res_S[\operatorname{tr}(m\adj M)/(f_1,f_2,h_1,h_2)]$; hence
$A\not\equiv0$ in the fifteen types of Table~\ref{tab:types} (all types with $d\le8$ and $K\ne0$, and
$(1,1)+(4,4)$), in each of which the hypothesis is proved for one of the two curves. Supports of lower status: exact computations in six types
(conditional on an identification without written proof), a Lean-certified finite computation at one sextic,
and computations modulo a prime. For a line $L$ and a disjoint $(3,3)$ curve $C$ on an octic---the classes of
Movasati's conjectural counterexample to Harris's conjecture---$A$ is a $2\times2$ matrix $\cD$ of residues,
linear in $M$. At every smooth octic containing such a pair and for every $r$ other than $0$ and $-1$, the only
component of the locus of $[L]+r[C]$ through that point of codimension smaller than $35$ is $\VZ$, the same for
all $r$; the others (present exactly where $\cD=0$, e.g.\ at the Fermat octic) have the general codimension
$35$. For $r=-1$ the locus is smooth of codimension $31$, a branch of the closure of the known family of octics
containing a complete intersection of two quartics. All of this is local at points of $\VZ$: nothing is claimed
about components of the Noether--Lefschetz locus that do not meet $\VZ$, nor about Harris's conjecture itself.
Movasati's statements in this direction are conjectural in his papers, as we read them, and his rigorous propositions are not
contradicted (the slice statements on the condition that his slices are spanned by monomials, as they appear to
be). Nothing here bears on the Hodge conjecture. The new parts were each read by one AI referee (Lemma~\ref{lem:branch} by its author
and one later referee); no human has checked anything.
\end{abstract}

\section{Introduction}\label{sec:intro}

\run{Setting.}
Let $S=\CC[x_0,\dots,x_3]$, with graded pieces $S_j$, and let $X=\{F=0\}\subset\PP^3$ be a smooth surface of
degree $d$; put $p_g=\dim S_{d-4}$ and let $H$ be the hyperplane class. We work in the space $T=S_d$ of forms of
degree $d$, with germs at $F$ (codimensions are the same in $\PP(S_d)$; tangent spaces are written as subspaces of
$S_d$). For a rational class $\delta\in H^2(X,\QQ)$ of type $(1,1)$, the \emph{Noether--Lefschetz locus}
$\NL(\delta)$ is the germ at $X$ of the set of surfaces $X'$ at which the flat transport of $\delta$ is still of
type $(1,1)$, with its natural structure of analytic space: the zero scheme of $p_g$ holomorphic functions
$\Phi_Q$, $Q\in S_{d-4}$ (fact (G1) of \S\ref{sec:facts}). A component of the Noether--Lefschetz locus of
surfaces of degree $d$ is called \emph{special} if its codimension is smaller than $p_g$, and \emph{general} if it
equals $p_g$.

\run{The proposal.}
A conjecture attributed to J.~Harris says that there are only finitely many special components. Voisin gave
counterexamples for large $d$ \cite{Vo91} and proved the finiteness of the \emph{reduced} special components for
$d=6,7$ \cite{Vo90}; we did not open these two papers and take the attributions from the abstract and the
reference list of \cite{Mo21}. Movasati \cite{Mo21} (abstract opened by us; body machine-read by one agent) presents evidence
that, for all but finitely many $r$, the Noether--Lefschetz loci of $C_1+rC_2$---$C_1$ a line and $C_2$ a
disjoint complete intersection of type $(3,3)$ on a smooth surface of degree~$8$---are distinct subvarieties of
codimension $31$ meeting in codimension $32$, the maximal codimension being $35$; the abstract calls this ``a
conjectural description of a counterexample'' to Harris's conjecture; the introduction says that the case $d=8$ is open. As
read by a third agent through machine extraction of the arXiv text: \cite{Mo21} contains no numbered conjecture
on these classes; it defines a \emph{special pencil} (tangent codimension below $p_g$ for all $r$, no inclusion
between the tangent spaces of the two loci, and smoothness of the locus as an analytic scheme for all but finitely
many $r$), states the expectation that the octic case provides one, and proves by computer (its Theorem~2) that at
the Fermat octic, on a $32$-parameter slice transversal to $\VZ$, the locus is ``$5$-smooth'' (\emph{$N$-smooth}, in his terminology as we read it: the $N$-th order
truncation of the locus is that of a smooth variety) for $r=r_2/r_1$, $1\le r_1\le10$, $|r_2|\le10$ (a range that contains $r=-1$); its Theorem~1
says that in the other cases it considers, for $4\le d\le8$, the locus is singular as an analytic scheme for the
coefficients in a stated range. The
preprint \cite{Mo26} (abstract opened; body machine-read by three of our agents) treats the same classes at the
Fermat octic. As read by referee~6: its Conjecture~1 says that for rational $r\neq0$ the local loci are
smooth, distinct analytic spaces of codimension $31$ through a locus of codimension $32$; its Proposition~3 says
that the locus is $2$-smooth in the full family; smoothness to higher order is obtained only on slices; and its
Remark~1 allows that the loci lie in a component of smaller codimension.

Non-reduced Noether--Lefschetz loci attached to sums of two curves with unequal coefficients are known: Maclean
\cite{Mac} (not opened) introduced a second-order invariant and treated quintics with two coplanar lines; Dan
\cite{Dan1} gave a criterion and examples for $d\ge5$, and \cite{Dan2} infinitely many non-reduced components of
non-maximal codimension for every $d\ge6$ (abstracts opened; per referee~7's reading, both papers use two
coplanar lines). For the distribution of exceptional components see also \cite{BKU} (search snippet only).
Kloosterman \cite[Thm~1.5, Thm~3.16]{Kl} treats two complete intersections $Y_1,Y_2$ of dimension $k$ meeting
in a complete intersection of codimension $c$ in each ($1\le c\le k+1$), on a hypersurface of degree $d$ in
$\PP^{2k+1}$. As read by the coordinating thread through a machine digest (to be checked against the paper; the index range
and the clause on the excess agree with a second machine reading, by referee~9, which is not a human reading): if
$\sum_{i=0}^{c-1}(e_i+d_i)<(c-1)d$, then for all but finitely many $\lambda$ there is no excess tangent dimension
and $\NL([Y_1]+\lambda[Y_2])=\NL([Y_1])\cap\NL([Y_2])$; for two disjoint curves on a surface ($k=1$, $c=2$) the
inequality reads $a_1+a_2+e_1+e_2<d$. If our reading of the index range of his inequality is right (it is an
inference from a machine reading and must be checked against the paper), Theorem~A concerns exactly the case of
equality, which his theorem excludes; there the excess is one. His paper \cite{Kl2} on linear subspaces meeting
in codimension one does not treat disjoint curves (reading of one agent). Duque Franco and Villaflor Loyola
\cite[\S5]{DFV} (read by the coordinating thread through machine digests of the arXiv HTML) treat pairs of
linear cycles meeting in a $\PP^m$ on the Fermat variety of degree $d$ and dimension $n$: the quadratic
fundamental form of the Hodge locus at the Fermat point vanishes for
$(d,m)=(3,\frac n2-3),(3,\frac n2-2),(4,\frac n2-2)$ (their Prop.~5.1--5.3) and is nonzero for $m=\frac n2-1$,
$d\ge2+\frac8n$ and unequal coefficients (their Thm~5.1). This concerns linear cycles, not disjoint curves.
In our first paper \cite{P1} the analogous question for two linear subspaces $P_1,P_2$ of a hypersurface was
treated: there the special coefficient is $\lambda_0=(-1)^{c-1}$, $c$ being the codimension of $P_1\cap P_2$ in
$P_i$; two disjoint lines have $c=2$ and $\lambda_0=-1$, consistent with the value $r=-1$ below.

\run{Results.}
A \emph{complete intersection of type $(a,b)$} is a subscheme of $\PP^3$ defined by a regular sequence of two
forms of degrees $a,b$; it may be singular, reducible or non-reduced.

\begin{thmA}
Let $Z_1=\{f_1=f_2=0\}$ and $Z_2=\{h_1=h_2=0\}$ be complete intersections in $\PP^3$ with $\deg f_i=a_i$,
$\deg h_j=e_j$ and $Z_1\cap Z_2=\emptyset$, and let $X=\{F=0\}$ be a smooth surface of degree
$d=a_1+a_2+e_1+e_2$ containing $Z_1$ and $Z_2$. Let $V_i$ be the germ at $X$ of the image of a neighbourhood of
$(Z_i,X)$ in the incidence variety of pairs \textup{(}complete intersection of the type of $Z_i$, surface
containing it\textup{)}, let $\VZ=V_1\cap V_2$, and for $r\in\QQ\setminus\{0\}$
let $\NL(r)=\NL([Z_1]+r[Z_2])$.
\begin{enumerate}[label=\textup{(\arabic*)}]
\item $F=f^{\mathsf T}Mh$ for a $2\times2$ matrix $M$ of forms, $\deg M_{ij}=d-a_i-e_j$; $\det M$ has degree $d$.
The germs $V_1,V_2$ are smooth at $X$, of codimensions $c_1,c_2$ \textup{(}Proposition~\ref{prop:first}\textup{)},
and meet transversally; $V_i=\NL([Z_i])$ as germs of analytic spaces; $\VZ$ is smooth of codimension $c_1+c_2$
and $T\VZ=T\,\NL([Z_1])\cap T\,\NL([Z_2])$.
\item For every $r\neq0$, $\codim T\,\NL(r)=c_1+c_2-1$: exactly one excess direction. The space
$K:=\{Q\in S_{d-4}:d\Phi_Q(X)=0\}$ does not depend on $r$, and $\dim K=p_g-(c_1+c_2-1)$.
\item Let $\cM_r\supset\NL(r)$ be the smooth germ of codimension $c_1+c_2-1$ cut out by $c_1+c_2-1$ of the
equations with independent differentials, and $\VZ=\{y=0\}\subset\cM_r$. Then
$\NL(r)=\{y^2u'_Q=0,\ Q\in K\}$ for holomorphic functions $u'_Q$ on $\cM_r$, where $\theta_Q:=\Phi_Q|_{\cM_r}
=y^2u'_Q$. There are forms $\gt_1,\gt_2\in S_d$, not depending on $r$, such that $G_r:=\gt_1-\gt_2/r$ lies in
$T\,\NL(r)\setminus T\VZ$, and a linear form $A\colon K\to\CC$ \emph{independent of $r$}, such that
\[
u'_Q(X)\,dy(G_r)^2=\tfrac12\Hess\theta_Q(G_r)=\frac{r+1}{r}\,A(Q)\qquad\text{for all }Q\in K,\ r\neq0 .
\]
\item $r=-1$: $\NL(-1)=\cM_{-1}$ is smooth and reduced of codimension $c_1+c_2-1$. It contains the surfaces
$\{F'+t\det M'=0\}$ for $X'=\{F'=f'^{\mathsf T}M'h'=0\}\in\VZ$ near $X$ and small $t$; for $t\neq0$ each of them
contains a complete intersection of type $(a_1+e_1,a_2+e_1)$ whose class is the transport of
$[Z_1]-[Z_2]+e_1H$, and $\NL(-1)$ is a branch of the closure of the locus of smooth surfaces containing such
a complete intersection \textup{(}Lemma~\ref{lem:branch}\textup{)}.
\item $r\notin\{0,-1\}$: let $W_r:=\{u'_Q=0\ \forall Q\in K\}\subset\cM_r$, so that
$\NL(r)_{\mathrm{red}}=\VZ\cup W_r$. The condition ``$A\neq0$ at $X$'' is intrinsic and is the same for all such
$r$. If $A\neq0$ at $X$, then $W_r=\emptyset$ and $\NL(r)=\{y^2=0\}$: it is $\VZ$ with a double structure, of
codimension $c_1+c_2$. If $A$ is not identically zero on the \textup{(}connected\textup{)} incidence variety
$\cI^\circ$ of triples $(Z_1,Z_2,X)$ as above, then at \emph{every} point of $\cI^\circ$ no component of $\NL(r)$
has codimension $c_1+c_2-1$, and $W_r\subsetneq\cM_r$ is a proper analytic subset.
\item \emph{A residue formula on $I(Z_1)I(Z_2)$.} Let $\kappa_1\ne0$ be as in Lemma~\ref{lem:pi}, let
$\varepsilon=\Res_S[\,\cdot\,/(f_1,f_2,h_1,h_2)]$, and let $A$ be normalised as in Lemma~\ref{lem:AB}. Suppose that
hypothesis \textup{(H)} of \S\ref{sec:F} holds for a general complete intersection of the type of $Z_1$, or for
a general one of the type of $Z_2$; it is proved if that curve is a line, a smooth conic, a smooth plane curve
of degree $a$ with $3a-6<2d-4-(a-1)$, or a smooth curve of type $(2,2)$. Then for all forms $m_{ij}$ of degree
$d-4-a_i-e_j$
\begin{equation}\label{eq:F}
A\Bigl(\sum_{i,j}f_ih_jm_{ij}\Bigr)=\frac1{2\kappa_1}\,\varepsilon\bigl(\operatorname{tr}(m\adj M)\bigr)
=\frac1{2\kappa_1}\sum_{i,j}\varepsilon\bigl(m_{ij}\adj(M)_{ji}\bigr).\tag{F}
\end{equation}
Consequently $A\not\equiv0$ on $\cI^\circ$ in every such type with $\max(a_1,a_2)+\max(e_1,e_2)\ge4$; this
covers all types of Table~\textup{\ref{tab:types}} with $K\ne0$. The formula describes $A$ on
$(I(Z_1)I(Z_2))_{d-4}\subset K$, which is all of $K$ if and only if all $a_i,e_j\le3$.
\end{enumerate}
\end{thmA}

Parts (1)--(5) do not decide whether $A\neq0$; this is a separate statement for each type. If $K=0$ (types
$(1,1)+(1,1)$, $(1,1)+(1,2)$, $(1,1)+(2,2)$) then $\NL(r)=\cM_r$ is smooth for every $r$ and there is nothing to
decide. In the $15$ types of Table~\ref{tab:types} with $K\ne0$ (the $14$ with $d\le8$, and $(1,1)+(4,4)$,
$d=10$) the statement $A\not\equiv0$ is \emph{proved} by Theorem~A(6): hypothesis (H) is proved by hand
for one of the two curves in each type (this part was read by one AI referee in a follow-up). Three further
supports, of lower status, are kept apart from it here and in Tables~\ref{tab:status}, \ref{tab:types}. They go
through the Jacobian-ring expression $\Ob$ for the second-order term (\S\ref{sec:comp}), whose identification
with that term is not written out as a proof (``the $\Ob$-link''): \emph{exact} values over $\QQ$ at one explicit
surface in the six types with $d\le7$ (for $(1,1)+(3,3)$ the exact residues of Theorem~B need no link); a
\emph{Lean-certified} finite computation at one explicit sextic with two disjoint conics, for every rational
$r\ne0$ (smoothness and disjointness are not certified); and one random point \emph{modulo one prime} per type,
not lifted to characteristic $0$. So in these types $\NL(r)$, $r\notin\{0,-1\}$, is $\VZ$ doubled at a general
point.

\begin{thmB}
Let $X=\{F=0\}$ be a smooth octic surface containing the line $L=\{x_0=x_1=0\}$ and a complete intersection
$C=\{h_1=h_2=0\}$ of two cubics with $L\cap C=\emptyset$. Write $f=(x_0,x_1)$, $F=f^{\mathsf T}Mh$, $g=Mh$, and on
$L$, with coordinates $(\sigma,\tau)=(x_2,x_3)$: $\bar h_j=h_j|_L$, $\Lambda=M|_L$,
$q_i=g_i|_L=(\Lambda\bar h)_i$ and $\varepsilon_L=\Res_{R_L}[\,\cdot\,/(\bar h_1,\bar h_2)]$, $R_L=S/I(L)=\CC[\sigma,\tau]$. Let
$\NL(r)=\NL([L]+r[C])$, $r\in\QQ\setminus\{0\}$, and use the notation of Theorem~A.
\begin{enumerate}[label=\textup{(\arabic*)}]
\item $T\,\NL([L])$, $T\,\NL([C])$, $T\VZ$, $T\,\NL(r)$ \textup{(}every $r\ne0$\textup{)} have codimensions
$5,27,32,31$, and $K=\langle f_ih_j\rangle=I(L\cup C)_4$ has dimension $4$.
\item For the normalised excess vector $G_r$ of \S\ref{sec:octic} and the constant $\kappa\neq0$ of
Lemma~\ref{lem:oct}\textup{(a)},
\[
u'_Q(X)\,dy(G_r)^2=\tfrac12\Hess\theta_Q(G_r)=\frac1{2\kappa}\cdot\frac{r+1}{r}\cdot\cD(Q),\qquad
\cD\Bigl(\sum a_{ij}f_ih_j\Bigr):=\sum a_{ij}\cD_{ij},
\]
\[
\cD_{1j}=\Res_L\Bigl[\frac{\bar h_j(\det\Lambda)^2\,d\sigma\wedge d\tau}{q_1^2,\ q_2}\Bigr],\qquad
\cD_{2j}=\Res_L\Bigl[\frac{\bar h_j(\det\Lambda)^2\,d\sigma\wedge d\tau}{q_1,\ q_2^2}\Bigr]
\quad\text{\textup{(}ordered denominators\textup{)}.}
\]
Moreover $\cD$ is linear in $M$: $\cD_{ij}=\varepsilon_L\bigl(\adj(\Lambda)_{ji}\bigr)$
\textup{(}Lemma~\ref{lem:L}\textup{)}. Hence $\cD(X)=0$ if and only if $M\equiv0$ modulo $(x_0,x_1,h_1,h_2)$, and
$\{\cD=0\}$ is irreducible of codimension exactly $4$ in the incidence variety $\cI^\circ$.
\item $r=-1$: $\NL(-1)=\cM_{-1}$ is smooth and reduced of codimension $31$; $\{F+t\det M=0\}$ lies in it for
small $t$ and, for $t\ne0$, contains a complete intersection $D_t$ of two quartics, of class the transport of
$[L]-[C]+3H$; $\cM_{-1}$ is a branch of the closure of the locus of octics containing a complete intersection
of two quartics \textup{(}Lemma~\ref{lem:branch}\textup{)}.
\item $r\notin\{0,-1\}$, at \emph{every} smooth $X$ as above. If $\cD(X)\neq0$, then
$\NL(r)=\{y^2=0\}\subset\cM_r$: $\VZ$ doubled, of codimension $32$, and nothing else. If $\cD(X)=0$, then
$\NL(r)=2\VZ\cup W_r$ near $X$ \textup{(}ideal $(y^2)\cap(u'_Q)_{Q\in K}$\textup{)}, where
$W_r=\{u'_Q=0,\ Q\in K\}$ is smooth of codimension exactly $35=p_g$, $W_r\cap\VZ=\{\cD=0\}$ is smooth of
codimension $36$, and each $\theta_Q$, $Q\in K\setminus\{0\}$, has zero $2$-jet and nonzero $3$-jet
$dy^2\cdot du'_Q$ on $\cM_r$. Hence $\VZ$ is the only component of $\NL(r)$ through $X$ of codimension smaller
than $35$, and it is the same for all $r$; in particular no component has codimension $31$. For $r\neq r'$ the
sets $W_r\setminus\VZ$ and $W_{r'}\setminus\VZ$ are disjoint.
\item At a split octic $F=P(x_0,x_1)+\tilde P(x_2,x_3)$, $P=\phi_1\phi_3\phi_4$,
$\tilde P=\gamma_1\gamma_3\gamma_4$ \textup{(}$\deg\phi_i=\deg\gamma_i=i$\textup{)}, with
$L=\{\phi_1=\gamma_1=0\}$ and $C=\{\phi_3=\gamma_3=0\}$---in particular at the Fermat octic with a line and a $(3,3)$ curve of this form---one has $\cD=0$, so the second case of \textup{(4)} holds there
\textup{(}\textup{(1)--(4)} being applied with $f=(\phi_1,\gamma_1)$ in place of $(x_0,x_1)$\textup{)}. Nothing
is claimed about $W_r$ beyond its germs at points of $\VZ$.
\end{enumerate}
\end{thmB}

Parts (1)--(3) are proved on paper. Part (4) follows from (2) and the linearity of $\cD$ in $M$, without computation (version~1 used an exact
value of $\cD$ and gave only codimensions between $32$ and $35$); (5) is a special case. The formula in (2) is the form $A$ of Theorem~A made explicit, $A=\cD/2\kappa$; by
Lemma~\ref{lem:L}, and because $\kappa=\kappa_1$ (\S\ref{sec:octic}), it is \eqref{eq:F} for this type.

\begin{cor*}
Let $X$ be any smooth octic surface containing a line $L$ and a disjoint complete intersection $C$ of two
cubics, and let $r\in\QQ\setminus\{0,-1\}$. Then the only component of the germ $\NL([L]+r[C])$ at $X$ of
codimension smaller than $35$ is $\VZ$, of codimension $32$, the same for all $r$; the other components are present exactly where $\cD=0$ and have the general codimension $35$.
In particular no component has codimension $31$, reduced or not. For $r=-1$ the germ is smooth and reduced of
codimension $31$: a branch of the closure of the known family of octics containing a complete intersection of
two quartics. All of this is local at points of $\VZ$: nothing is claimed about components of the
Noether--Lefschetz locus of octics that do not meet $\VZ$, nor about Harris's conjecture itself.
\end{cor*}

\run{A low-degree instance.}
The first cases with $K\ne0$ occur for $d=6$: a line and a disjoint plane cubic ($\dim K=2$), and two disjoint
conics $C_1,C_2$ ($\dim K=1$, $K=\CC\cdot\ell\ell'$ with $\ell,\ell'$ the planes of the conics). For two disjoint
conics on a sextic, $\VZ$ has codimension $10=p_g$ and $T\,\NL(r)$ has codimension $9$; $\NL(C_1-C_2)$ is smooth
of codimension $9$, a branch of the closure of the locus of sextics containing a complete intersection of a
quadric and a cubic, of class $C_1-C_2+H$ (proved); and for $r\notin\{0,-1\}$, $\NL(C_1+rC_2)$ is $\VZ$ with a
double structure at a general point---a non-reduced locus of the expected codimension---because
$A\not\equiv0$. For this type \eqref{eq:F} reads $A(\ell\ell')=\varepsilon(M_{22})/2\kappa_1$, $M_{22}$ being the quadric entry of $M$; for the other supports see \S\ref{sec:comp}. This
instance appears to be new: we found no earlier source in a finite search.

\run{What is and is not claimed.}
Known, and to be attributed: non-reduced Noether--Lefschetz loci coming from sums of two curves with unequal
coefficients (Maclean \cite{Mac}; Dan \cite{Dan1,Dan2}); the second-order invariant (Maclean \cite{Mac}) (as attributed
in \cite{Mo21}; \cite{Mac} was not opened);
quadratic fundamental forms of Hodge loci of pairs of linear cycles at Fermat points \cite{DFV}; and the family
of surfaces containing a complete-intersection curve as a special component---so $\NL(-1)$ is a known locus, and
nothing here produces a special component that was not known before. What appears to be new (no earlier source was found in finite searches; see \S\ref{sec:open} for what was not read) is: the disjoint complete-intersection configuration with
$d=a_1+a_2+e_1+e_2$ (the case of equality in \cite[Thm~1.5]{Kl}, as we read it); the law $(r+1)/r$; the
identification of $r=-1$ with the determinantal family; the residue formula \eqref{eq:F}; and the conclusion
that, at the points of $\VZ$, Movasati's classes give no component of codimension smaller than $35$ other than
$\VZ$ and, for $r=-1$, a known locus. This concerns germs at points of $\VZ$: the proposed
mechanism (distinct loci of codimension $31$ through $\VZ$) does not occur; on Harris's conjecture for octic
surfaces nothing is said. Movasati's statements in that direction are conjectural in his papers, and his rigorous propositions, as
we read them, are not contradicted---the slice statements on the condition that his slices are spanned by
monomials (Remark~\ref{rem:sym}, Table~\ref{tab:mov}). Nothing here bears on the Hodge conjecture: all classes
considered are classes of curves.

\begin{table}[!t]
\centering\footnotesize
\begin{tabular}{@{}p{0.33\textwidth}p{0.42\textwidth}p{0.20\textwidth}@{}}
\toprule
Statement & Method & Status\\
\midrule
\rg Facts (G1), (G2), (A), (R), (HB), (T), (CI); period formula for complete-intersection curves &
 \rg \cite{Gr,CG,Ei,VL}, none opened (\cite{VL} machine-read); no source named for (A), (R), (T), (CI) &
 \rg Standard; cited, or no source given\tabularnewline
\rg Theorems A(1)--(5), B(1)--(3), except the last clause of A(4) and of B(3) (next row) & \rg \S\S\ref{sec:first}--\ref{sec:cons}, \ref{sec:octic}; one AI referee
 each (``referee~7'', ``referee~6''); re-derived by referee~8 & \rg Proved on paper, AI-refereed\tabularnewline
\rg Lemma~\ref{lem:branch} & \rg proof supplied by referee~8; read by one later AI referee (referee~9) only &
 \rg Proved on paper; read by one later AI referee\tabularnewline
\rg Theorem A(6): \eqref{eq:F}, Lemma~\ref{lem:R}; hence $A\not\equiv0$ in the $15$ types of
 Table~\ref{tab:types} & \rg \S\ref{sec:F}; (H) by hand for one curve in each type; working notes read by one AI
 referee in a follow-up (``referee~7b'') & \rg Proved on paper; AI-refereed (follow-up)\tabularnewline
\rg Hypothesis (H) for types without a line, a conic, a plane curve in the range or a $(2,2)$ curve &
 \rg $t_0$ of random curves modulo $p$ & \rg Observed modulo $p$\tabularnewline
\rg Lemma~\ref{lem:L}, Corollary~\ref{cor:L}, Theorem B(4), (5) & \rg \S\S\ref{ss:every}, \ref{ss:split}; working notes read by one
 AI referee in a follow-up (``referee~6b''); exact checks at $6+11$ data & \rg Proved on paper; AI-refereed
 (follow-up)\tabularnewline
\rg Identities of Lemma~\ref{lem:lean}, over any commutative ring & \rg \path{lean/OcticDetCert.lean}, Part~A &
 \rg Lean-certified identity\tabularnewline
\rg $\cD(X_*)=\bigl(\begin{smallmatrix}0&-1\\-1&0\end{smallmatrix}\bigr)$ &
 \rg exact residues; reproduced by referees~6 and~8 & Exact over $\QQ$\\
\rg $A\not\equiv0$ in the types marked ``exact ($\Ob$)'' in Table~\ref{tab:types} &
 \rg exact evaluation of $\Ob$ at one explicit surface & \rg Exact over $\QQ$, conditional on the
 $\Ob$-link\tabularnewline
\rg $\Ob(r)=A_0(1+\frac1r)$, $A_0\ne0$, for all rational $r\ne0$, at one explicit sextic &
 \rg \path{lean/SexticCert.lean}; not certified: smoothness, disjointness, $\dim(S/J)_{16}=1$, the $\Ob$-link &
 \rg Lean-certified finite computation\tabularnewline
\rg $A\neq0$ at a random point ($15$ types); the law $(r+1)/r$ ($5$ types); \eqref{eq:F} ($14$ types); split
 pairs on Fermat surfaces & \rg \path{second_order.py}, \path{split_fermat.py}: one point, one prime, through
 $\Ob$; the referees' own programs & \rg Modulo $p$; conditional on the $\Ob$-link; not lifted to
 characteristic $0$\tabularnewline
\rg Fermat octic: direct third-order computation; the tangent codimensions $(1,9,9,10)$ of \cite{Mo26} on his $25$-dimensional family &
 \rg \path{octic/fermat3.py}, two primes, one implementation, examined by no referee & Modulo $p$\\
\rg Statements of \cite{Mo21,Mo26,Kl,Kl2,DFV,Mac,Dan1,Dan2,Vo90,Vo91,BKU} & \rg abstracts, snippets or machine
 extraction & \rg Cited, as read\tabularnewline
\bottomrule
\end{tabular}
\caption{Status of the statements used or reported. No human has checked any of them.}\label{tab:status}
\end{table}

\section{Standard facts used}\label{sec:facts}

The following facts are cited in the form in which they are used; they are not re-proved here, and none of the
sources \cite{Gr,CG,Ei} was opened for this paper. For (A), (R), (T), (CI) and for the scheme structure in (G1) we
name no source: they are standard, and are stated in the form used. $\Omega$ is the standard $3$-form on $\PP^3$.
\begin{itemize}
\item[(G1)] (Griffiths) For $N\in S_{ld-4}$ the form $N\Omega/F^l$ has a residue in $H^2_{\mathrm{prim}}(X)$;
$F^2H^2=\{\Res(Q\Omega/F):Q\in S_{d-4}\}$ and $F^1H^2_{\mathrm{prim}}=\{\Res(N\Omega/F^2):N\in S_{2d-4}\}$.
For a flat rational class $\delta$, of type $(1,1)$ at $X$, the locus $\NL(\delta)$ is the zero scheme of the
holomorphic functions $\Phi^\delta_Q(F')=\int_\delta\Res(Q\Omega/F')$, $Q\in S_{d-4}$ ($F'$ near $F$). These
depend on $\delta$ only through its primitive part (the tube map kills $H$), so $\NL(\delta)=\NL(\delta+cH)$.
\item[(G2)] (Griffiths--Dwork) For $m\in S_{ld-3}$: $m\,\partial_jF\,\Omega/F^{l+1}\equiv\frac1l\,\partial_jm\,
\Omega/F^l$ in $H^3(\PP^3\setminus X)$; and algebraic classes are orthogonal to $F^2H^2$.
\item[(A)] Let $\mathcal X\to B$ be a smooth projective family of surfaces and $\mathcal Z\subset\mathcal X$ a
closed subscheme, flat over $B$, whose fibres are curves without embedded points. Then $\mathcal Z$ is a relative
effective Cartier divisor and $t\mapsto[\mathcal Z_t]$, the first Chern class of $\mathcal O(\mathcal Z)$
restricted to $X_t$, is a flat section;
$[\mathcal Z_t]$ is the cycle class with multiplicities. In particular the transported class is of type $(1,1)$.
\item[(R)] (Grothendieck residues) For a regular sequence $\varphi=(\varphi_1,\dots,\varphi_4)$ of forms in $S$,
$S/(\varphi)$ is Gorenstein with socle in degree $\sigma=\sum\deg\varphi_i-4$; the functional
$\Res_S[\,\cdot\,/\varphi]\colon S_\sigma\to\CC$ is nonzero with kernel $(\varphi)_\sigma$; the multiplication
pairings $(S/\varphi)_j\times(S/\varphi)_{\sigma-j}\to(S/\varphi)_\sigma$ are perfect; and if $\varphi=P\psi$ for
a matrix $P$ of forms then $\Res[m/\psi]=\Res[m\det P/\varphi]$. The same holds for $n$ forms in $n$ variables; we use $n=4$ and $n=2$. The residue is alternating in the denominators, which are always ordered.
\item[(HB)] (Hilbert--Burch) If the ideal of $2\times2$ minors of a $2\times3$ matrix of forms has codimension
$2$, it is perfect, with graded resolution (hence Hilbert polynomial) determined by the degrees.
\item[(T)] $\operatorname{Tor}_1^S(S/I,S/I')=(I\cap I')/(I\cdot I')$.
\item[(CI)] The complete intersections of a given type form a smooth family $\cH$, with
$T_Z\cH=H^0(N_Z)=(S/I(Z))_{a_1}\oplus(S/I(Z))_{a_2}$, and $H^0(\mathcal O_Z(j))=(S/I(Z))_j$.
\end{itemize}
Complete-intersection curves, and the determinantal curves of \S\ref{sec:det}, are Cohen--Macaulay, so (A)
applies to them even when they are singular, reducible or non-reduced.

\run{Period functionals.} For a curve $Z_i\subset X$ put $\pi_i(N):=\int_{Z_i}\Res(N\Omega/F^2)$ for
$N\in S_{2d-4}$ and $\varphi_i(N'):=\int_{Z_i}\Res(N'\Omega/F^3)$ for $N'\in S_{3d-4}$; for
$\delta=\delta_r=[Z_1]+r[Z_2]$ put $\pi_\delta=\pi_1+r\pi_2$, $\varphi_\delta=\varphi_1+r\varphi_2$. The function
$\Phi^\delta_Q$ is the integral of $Q\Omega/F'$ over a fixed tube cycle in $\PP^3\setminus X$, and expanding
$1/(F+G)$ termwise gives
\begin{equation}\label{eq:exp}
\Phi^\delta_Q(F+G)=\Phi^\delta_Q(F)-\pi_\delta(QG)+\varphi_\delta(QG^2)-\cdots .
\end{equation}

\section{Algebra and first order}\label{sec:first}

Throughout \S\S\ref{sec:first}--\ref{sec:cons}, $Z_1,Z_2,X,F$ are as in Theorem~A.

\begin{lemma}\label{lem:alg}
\begin{enumerate}[label=\textup{(\alph*)}]
\item $(f_1,f_2,h_1,h_2)$ is a regular sequence, $I(Z_1)\cap I(Z_2)=I(Z_1)\cdot I(Z_2)$, and $F=f^{\mathsf T}Mh$
with $\deg M_{ij}=d-a_i-e_j>0$. Put $g:=Mh$ and $k:=M^{\mathsf T}f$, so that $F=f_1g_1+f_2g_2=h_1k_1+h_2k_2$.
\item $(f_1,f_2,g_1,g_2)$ and $(h_1,h_2,k_1,k_2)$ are regular sequences.
\item $(f,g)\subset(f,h)\supset(h,k)$.
\item $S/(f,h)$ has socle degree $d-4$. Hence $S_d=I(Z_1)_d+I(Z_2)_d$, and $(f,h)_{d-4}$ is a hyperplane in
$S_{d-4}$: the kernel of $\varepsilon:=\Res_S[\,\cdot\,/(f_1,f_2,h_1,h_2)]$.
\item $\Res_S[m\det M/(f,g)]=\varepsilon(m)=\Res_S[m\det M/(h,k)]$ for $m\in S_{d-4}$.
\end{enumerate}
\end{lemma}
\begin{proof}
(a) The sequence is regular because $Z_1\cap Z_2=\emptyset$. So $f_1,f_2$ is regular on $S/I(Z_2)$, the Koszul
complex gives $\operatorname{Tor}_1(S/I(Z_1),S/I(Z_2))=0$, and (T) gives the equality of ideals; hence
$F\in I(Z_1)I(Z_2)$. (b) At a common zero of $f,g$ one has $\nabla F=\sum g_i\nabla f_i+\sum f_i\nabla g_i=0$,
impossible on the smooth $X$. (c) $g=Mh$ and $k=M^{\mathsf T}f$. (d) The socle degree is
$a_1+a_2+e_1+e_2-4$ by (R), so $(S/(f,h))_j=0$ for $j>d-4$. (e) Transformation law (R) for
$(f_1,f_2,g_1,g_2)^{\mathsf T}=\operatorname{diag}(1,1,M)\,(f_1,f_2,h_1,h_2)^{\mathsf T}$; the socle degree of
$S/(f,g)$ is $2d-4$. For $(h,k)$ use $\operatorname{diag}(1,1,M^{\mathsf T})$ and the even permutation
$(h_1,h_2,f_1,f_2)$.
\end{proof}

The matrix $M$ is not unique: the admissible changes are
$M'=(f_2,-f_1)^{\mathsf T}(s_1,s_2)+(w_1,w_2)^{\mathsf T}(h_2,-h_1)$. Then $g$ changes by multiples of the $f_i$
and $k$ by multiples of the $h_j$, so $(f,g)$ and $(h,k)$ are unchanged; and
$\det(M+M')-\det M=s_1k_2-s_2k_1+w_1g_2-w_2g_1-(w\cdot f)(s\cdot h)$ lies in $(f,g)_d\cap(h,k)_d$, which is $T\VZ$
(Proposition~\ref{prop:first}). So nothing below depends on the choice.

\begin{lemma}\label{lem:pi}
$\pi_1=\kappa_1\Res_S[\,\cdot\,/(f,g)]$ and $\pi_2=\kappa_2\Res_S[\,\cdot\,/(h,k)]$ with $\kappa_1,\kappa_2\ne0$.
\end{lemma}
\begin{proof}
(i) If $G\in I(Z_1)_d$ then $Z_1\subset\{F+tG=0\}$, so $\Phi^{[Z_1]}_Q(F+tG)\equiv0$ by (A) and (G1), and
$\pi_1(QG)=0$ for all $Q\in S_{d-4}$ by \eqref{eq:exp}. (ii) If $G=g_1u$, $u\in S_{a_1}$, then
$F+tG=(f_1+tu)g_1+f_2g_2$ contains the complete intersection $\{f_1+tu=f_2=0\}$, a flat deformation of $Z_1$;
again $\pi_1(QG)=0$. The same holds for $g_2u$, $u\in S_{a_2}$. (iii) $S_{d-4}\cdot(f,g)_d=(f,g)_{2d-4}$, the
generators having degrees $\le d$. So $\pi_1$ kills $(f,g)_{2d-4}$, a hyperplane in $S_{2d-4}$ by (R).
(iv) $\pi_1\neq0$: by (G1), $[Z_1]_{\mathrm{prim}}=\Res(N_0\Omega/F^2)$ for some $N_0$, and
$\pi_1(N_0)=\langle Z_1,Z_1\rangle_{\mathrm{prim}}=a_1a_2(a_1+a_2-d)-(a_1a_2)^2/d<0$.
\end{proof}

\begin{proposition}[Theorem A(1)]\label{prop:first}
\begin{enumerate}[label=\textup{(\alph*)}]
\item $T\,\NL([Z_1])=(f,g)_d$ and $T\,\NL([Z_2])=(h,k)_d$, of codimensions
$c_1:=\dim(S/(f,g))_d=\dim(S/(f,g))_{d-4}$ and $c_2:=\dim(S/(h,k))_d=\dim(S/(h,k))_{d-4}$.
\item $V_1$ is smooth at $X$ with $TV_1=(f,g)_d$, and $V_1=\NL([Z_1])$ as germs of analytic spaces; likewise
for $V_2$.
\item $V_1$ and $V_2$ are transversal; $\VZ$ is smooth of codimension $c_1+c_2$, and $T\VZ=(f,g)_d\cap(h,k)_d$.
\end{enumerate}
\end{proposition}
\begin{proof}
(a) By \eqref{eq:exp}, $T\,\NL([Z_1])=\{G:\pi_1(QG)=0\ \forall Q\in S_{d-4}\}$; this is $(f,g)_d$ by
Lemma~\ref{lem:pi} and the perfect pairing (R) in degrees $(d-4,d)$.
(b) The incidence variety $\cI_1=\{(Z,X'):Z\subset X'\}$ is a projective bundle over $\cH$ (CI), hence smooth. A
tangent vector at $(Z_1,X)$ is a pair $((u_1,u_2),G)$ with $G\equiv u_1g_1+u_2g_2$ mod $I(Z_1)$ (for
$Z_t=\{f_1+tu_1=f_2+tu_2=0\}$, as in the proof of Lemma~\ref{lem:pi}(ii)). So the image of $T\cI_1$ in $S_d$ is
$(f,g)_d$, and the kernel is
$\{(u_1,u_2)\in(S/I(Z_1))_{a_1}\oplus(S/I(Z_1))_{a_2}:u_1g_1+u_2g_2\in I(Z_1)\}$. As $g_1,g_2$ is a regular
sequence on $S/I(Z_1)$ (Lemma~\ref{lem:alg}(b)), $(u_1,u_2)=s\cdot(g_2,-g_1)$ with $\deg s=a_1+a_2-d<0$: the
kernel is $0$. So $V_1$ is smooth with the stated tangent space. $V_1\subset\NL([Z_1])$ by (A), and the tangent
spaces agree by (a); a germ of analytic space containing a smooth germ of dimension equal to its embedding
dimension coincides with it.
(c) $TV_1+TV_2\supset I(Z_1)_d+I(Z_2)_d=S_d$ by Lemma~\ref{lem:alg}(d).
\end{proof}

\begin{proposition}[Theorem A(2)]\label{prop:excess}
For every $r\ne0$, $T\,\NL(r)/T\VZ$ is a line, so $\codim T\,\NL(r)=c_1+c_2-1$. Moreover
$K:=\{Q:d\Phi_Q(X)=0\}=(f,g)_{d-4}\cap(h,k)_{d-4}$, independently of $r$, and $\dim K=p_g-(c_1+c_2-1)$.
\end{proposition}
\begin{proof}
Put $u_i(G):=\pi_i(\,\cdot\,G)\in S_{d-4}^*$. By (R), $u_1$ maps $S_d$ onto $U_1:=((f,g)_{d-4})^\perp$ and $u_2$
onto $U_2:=((h,k)_{d-4})^\perp$, and $\dim U_i=c_i$. By Proposition~\ref{prop:first}(c) the map
$(u_1,u_2)\colon S_d\to U_1\oplus U_2$ has kernel $T\VZ$, of codimension $c_1+c_2$, so it is onto. Now
$T\,\NL(r)=\{u_1+ru_2=0\}$, so $T\,\NL(r)/T\VZ\cong U_1\cap U_2=((f,g)_{d-4}+(h,k)_{d-4})^\perp$. By
Lemma~\ref{lem:alg}(c) the sum lies in $(f,h)_{d-4}$, and it contains $(f)_{d-4}+(h)_{d-4}=(f,h)_{d-4}$. So
$U_1\cap U_2=\CC\varepsilon$ is a line, for every $r\ne0$. Finally
$\{Q:d\Phi_Q(X)=0\}=(U_1+U_2)^\perp=(f,g)_{d-4}\cap(h,k)_{d-4}$.
\end{proof}

\section{Structure of \texorpdfstring{$\NL(r)$}{NL(r)} and the dependence on \texorpdfstring{$r$}{r}}\label{sec:struct}

Fix $r\neq0$ and put $n=c_1+c_2-1$. Choose $Q'_1,\dots,Q'_n\in S_{d-4}$ with independent differentials
$d\Phi_{Q'_i}(X)$ and let $\cM_r:=\{\Phi_{Q'_i}=0\}$, a smooth germ. Then $S_{d-4}=K\oplus\langle Q'\rangle$ and
$\NL(r)=\{\theta_Q=0,\ Q\in K\}$ with $\theta_Q:=\Phi_Q|_{\cM_r}$.

\begin{lemma}\label{lem:struct}
$\VZ$ is a smooth hypersurface $\{y=0\}$ of $\cM_r$, and $\theta_Q=y^2u'_Q$ with $u'_Q$ holomorphic on $\cM_r$.
For $Q\in K$ and $G\in T_X\cM_r=T\,\NL(r)$,
\begin{equation}\label{eq:hess}
u'_Q(X)\,dy(G)^2=\tfrac12\Hess\theta_Q(G)=\varphi_\delta(QG^2),
\end{equation}
and $T\VZ$ lies in the kernel of this quadratic form.
\end{lemma}
\begin{proof}
$\VZ\subset\NL(r)\subset\cM_r$ and $\VZ$ is smooth of codimension $n+1$, so $\theta_Q=y\,u_Q$. Every $X'\in\VZ$
near $X$ is again a smooth surface with a disjoint pair of the same types, and the transport of $\delta_r$ is
$[Z_1']+r[Z_2']$ by (A); so \S\ref{sec:first} holds at $X'$: $\codim T_{X'}\NL(r)=n$, hence
$T_{X'}\NL(r)=T_{X'}\cM_r$ and $d\theta_Q(X')=u_Q(X')\,dy=0$. So $u_Q=y\,u'_Q$. Since $d\Phi_Q(X)=0$ for
$Q\in K$, the Hessian of the restriction to $\cM_r$ is the restriction of the Hessian, which is
$2\varphi_\delta(QG^2)$ by \eqref{eq:exp}.
\end{proof}

At $X'\in\VZ$ the same holds with $Q$ replaced by its projection $Q_{X'}\in K_{X'}$ along $\langle Q'\rangle$,
since $\Phi_{Q_{X'}}|_{\cM_r}=\theta_Q$.

\begin{lemma}\label{lem:AB}
There are $\gt_1\in I(Z_2)_d$ with $u_1(\gt_1)=\varepsilon$ and $\gt_2\in I(Z_1)_d$ with
$u_2(\gt_2)=\varepsilon$. For such a choice, $G_r:=\gt_1-\gt_2/r$ lies in $T\,\NL(r)\setminus T\VZ$ for every
$r\ne0$, and for $Q\in K$
\[
\varphi_\delta(QG_r^2)=A(Q)+\frac{B(Q)}{r},\qquad
A(Q):=\varphi_1(Q\gt_1^2)-2\varphi_2(Q\gt_1\gt_2),\quad B(Q):=\varphi_2(Q\gt_2^2)-2\varphi_1(Q\gt_1\gt_2),
\]
where $A$ and $B$ do not depend on $r$, nor on the choice of $\gt_1,\gt_2$.
\end{lemma}
\begin{proof}
$\varepsilon\in U_1=u_1(S_d)$, and $u_1(S_d)=u_1(I(Z_2)_d)$ because $S_d=I(Z_2)_d+I(Z_1)_d$ and $u_1$ kills
$I(Z_1)_d$; likewise for $\gt_2$. Then $u_2(\gt_1)=0=u_1(\gt_2)$, so $u_1(G_r)+ru_2(G_r)=\varepsilon-\varepsilon=0$,
and $G_r\notin T\VZ$ since $u_1(G_r)\ne0$. Next, $Z_1\subset\{F+s\gt_2=0\}$ for all $s$, so
$\Phi^{[Z_1]}_Q(F+s\gt_2)\equiv0$ and the $s^2$-coefficient in \eqref{eq:exp} gives $\varphi_1(Q\gt_2^2)=0$ for
every $Q\in S_{d-4}$; likewise $\varphi_2(Q\gt_1^2)=0$. Expanding
$(\varphi_1+r\varphi_2)\bigl(Q(\gt_1-\gt_2/r)^2\bigr)$, the terms with $1/r^2$ and with $r$ vanish, and what
remains is $A(Q)+B(Q)/r$. Finally, $\gt_1$ is unique modulo $I(Z_2)_d\cap(f,g)_d\subset T\VZ$, and likewise
$\gt_2$; another choice changes $G_r$ by an element of $T\VZ$, which is in the kernel of the form
\eqref{eq:hess}; so $A+B/r$ is unchanged for every $r$, and hence so are $A$ and $B$.
\end{proof}

\begin{proposition}[Theorem A(3)]\label{prop:law}
$B=A$. Hence $\tfrac12\Hess\theta_Q(G_r)=\frac{r+1}{r}A(Q)$ for all $Q\in K$ and $r\neq0$.
\end{proposition}
\begin{proof}
By Lemmas~\ref{lem:arc} and \ref{lem:detM} below (which do not use this section) there is a holomorphic arc
$X_t$ in $\NL(-1)$ with $X_0=X$ and tangent vector $\det M\notin T\VZ$. Take $r=-1$. Along the arc
$\theta_Q=y^2u'_Q$ vanishes and $y(X_t)$ has a simple zero at $t=0$, so $u'_Q(X_t)=0$ for $t\ne0$, hence
$u'_Q(X)=0$. By \eqref{eq:hess} and Lemma~\ref{lem:AB}, $0=\varphi_\delta(QG_{-1}^2)=A(Q)-B(Q)$.
\end{proof}

\begin{remark}\label{rem:cross}
When $K\subset I(Z_1)\cap I(Z_2)$ the cross terms vanish separately, $A(Q)=\varphi_1(Q\gt_1^2)$ and
$B(Q)=\varphi_2(Q\gt_2^2)$; this is the situation of \S\ref{sec:octic}. In general $K$ is bigger: for
$(1,1)+(1,4)$, $d=7$, the cubic $k_2$ lies in $K$. Only the shape $A+B/r$ of the expansion is used, not its
coefficients.
\end{remark}

\section{The determinantal family at \texorpdfstring{$r=-1$}{r = -1}}\label{sec:det}

\begin{lemma}[ring identities; Lean-certified]\label{lem:lean}
Let $f=(f_1,f_2)^{\mathsf T}$, $h=(h_1,h_2)^{\mathsf T}$, $M=(M_{ij})$ and $t$ be elements of a commutative ring;
put $F=f^{\mathsf T}Mh$, $k_2=M_{12}f_1+M_{22}f_2$, $F_t=F+t\det M$, $N_t=t\adj M+hf^{\mathsf T}$, with first row
$(n_{11},n_{12})=(f_1h_1+tM_{22},\ f_2h_1-tM_{12})$, and
$H_t=\bigl(\begin{smallmatrix}f_1&f_2&t\\-M_{22}&M_{12}&h_1\end{smallmatrix}\bigr)$. Then:
\begin{enumerate}[label=\textup{(\alph*)}]
\item $\det N_t=t^2\det M+t\,f^{\mathsf T}Mh=t\,F_t$;
\item the $2\times2$ minors of $H_t$ are $k_2,\ n_{11},\ n_{12}$, and $t\,k_2=f_2n_{11}-f_1n_{12}$,
$h_1k_2=M_{12}n_{11}+M_{22}n_{12}$;
\item $F_t=M_{11}n_{11}+M_{21}n_{12}+h_2k_2$, and
$t\,F_t=(tM_{11}+h_2f_2)\,n_{11}+(tM_{21}-h_2f_1)\,n_{12}$;
\item $F=M_{11}\cdot f_1h_1+M_{21}\cdot f_2h_1+h_2k_2$.
\end{enumerate}
\end{lemma}
These six identities (theorems \texttt{det\_identity}, \texttt{hb\_minors}, \texttt{hb\_syzygies},
\texttt{family}, \texttt{family\_in\_ci}, \texttt{limit} of \path{lean/OcticDetCert.lean}, Part~A) are checked
by the Lean~4 kernel (core only, no Mathlib, no \texttt{native\_decide}, no \texttt{sorry}) over an arbitrary
commutative ring. Nothing else about this family is Lean-certified (the only other Lean-certified item
is the finite computation at one sextic, \S\ref{sec:comp}): not homogeneity, codimension, flatness, the
cycle classes, or any statement about Noether--Lefschetz loci.

In our graded situation $n_{11},n_{12}$ are homogeneous of degrees $a_1+e_1$, $a_2+e_1$. Let
$X_t=\{F_t=0\}$, smooth for small $t$, and let $D_t\subset\PP^3$ be the scheme of the minors of $H_t$. By (b),
(c): $D_t=\{n_{11}=n_{12}=0\}$ for $t\neq0$, and $D_t\subset X_t$ for every $t$, $t=0$ included.

\begin{lemma}[the fibre at $t=0$]\label{lem:Y}
Let $Y=D_0$, defined by $(k_2,f_1h_1,f_2h_1)$, and $C':=\{h_1=k_2=0\}$. Then $C'$ is a curve, $Y$ has support
$Z_1\cup C'$, $[Y]=[Z_1]+[C']$ as cycles, $Y\subset X$, and $[C']=e_1H-[Z_2]$ on $X$.
\end{lemma}
\begin{proof}
$h_1$ and $k_2$ have no common factor: it would divide $F=h_1k_1+h_2k_2$, and $X$ is irreducible of larger
degree. No component of $Z_1$ lies in $\{h_1=0\}$ (it would meet $\{h_2=0\}$, but $Z_1\cap Z_2=\emptyset$). So at
the generic points of $Z_1$ the ideal of $Y$ is $(f_1,f_2)$ ($h_1$ is a unit and $k_2\in(f)$), and off $Z_1$ it
is $(h_1,k_2)$; $Z_1$ and $C'$ may share points but no component. $Y\subset X$ by Lemma~\ref{lem:lean}(d). On
$X$, $F\equiv h_2k_2$ mod $h_1$, so $\operatorname{div}_X(h_1)=[Z_2]+[C']$.
\end{proof}

\begin{lemma}[the arc]\label{lem:arc}
The family $(D_t)$ is flat near $t=0$, the fibres are Cohen--Macaulay curves with the Hilbert polynomial of a
complete intersection of type $(a_1+e_1,a_2+e_1)$, and the transport of $[Z_1]-[Z_2]+e_1H$ to $X_t$ is the class
of $D_t$. In particular $X_t\in\NL(-1)$ for small $t$.
\end{lemma}
\begin{proof}
Let $I\subset S[t]$ be the ideal of the minors of $H_t$. Near $t=0$ its zero set in $\CC^4\times\CC$ has
dimension $3$ (the fibre over $0$ has dimension $2$ by Lemma~\ref{lem:Y}; semicontinuity). So, after inverting a
polynomial in $t$ with nonzero value at $0$, $I$ has codimension $2$, $S[t]/I$ is Cohen--Macaulay (HB), no
associated prime contains $t$ (the fibre over $0$ has dimension $2$), and $S[t]/I$ is torsion-free, hence flat,
over $\CC[t]$ near $0$. The fibres have no embedded points and $D_t\subset X_t$, so the family is a relative effective
Cartier divisor in the smooth family $(X_t)$ and (A) applies: the class of $D_t$ is the transport of
$[Y]=[Z_1]+e_1H-[Z_2]$. Since $\NL(\delta)$ depends only on the primitive part of $\delta$ (G1), $X_t\in\NL(-1)$.
\end{proof}

\begin{lemma}\label{lem:detM}
$\det M\notin(f,g)_d$; in particular $\det M\notin T\VZ$. Moreover $\kappa_1=\kappa_2$ in Lemma~\ref{lem:pi}.
\end{lemma}
\begin{proof}
By Lemma~\ref{lem:alg}(e) and Lemma~\ref{lem:pi}, $\pi_1(m\det M)=\kappa_1\varepsilon(m)$ and
$\pi_2(m\det M)=\kappa_2\varepsilon(m)$ for $m\in S_{d-4}$, and $\varepsilon\ne0$; this gives the first statement.
By Lemma~\ref{lem:arc}, $\det M$ is tangent to $\NL(-1)$, and
$T\,\NL(-1)=\{G:\pi_1(QG)=\pi_2(QG)\ \forall Q\in S_{d-4}\}$ by \eqref{eq:exp}; so $\kappa_1=\kappa_2$.
\end{proof}

\begin{proposition}[Theorem A(4)]\label{prop:minus1}
$\NL(-1)=\cM_{-1}$ as analytic spaces, and it contains the surfaces described in Theorem~\textup{A(4)}.
\end{proposition}
\begin{proof}
For $X'\in\VZ$ near $X$ choose $M(X')$ holomorphically (the curves move holomorphically with $X'$ by
Proposition~\ref{prop:first}, and $M$ solves a linear system of constant rank: $\dim(I(Z_1')\cap I(Z_2'))_d$ is
constant because $(S/(I(Z_1')+I(Z_2')))_d=0$). Over a connected neighbourhood $U$ of $(X,0)$ in $\VZ\times\CC$
the minors of $H(X',t)$ define a subscheme of $\PP^3\times U$ with fibres of dimension $1$ (semicontinuity from
$(X,0)$); as in Lemma~\ref{lem:arc}, the Hilbert--Burch complex over $\mathcal O_U[x_0,\dots,x_3]$ is exact and
stays exact on every fibre, so the subscheme is flat over $U$; after shrinking $U$ all surfaces
$\Psi(X',t):=F'+t\det M(X')$ are smooth. By (A) the class of the fibres is a flat section over the connected
$U$, and at $(X,0)$ it is $[Z_1]-[Z_2]+e_1H$ (Lemma~\ref{lem:Y}). Hence $\Psi(U)\subset\NL(-1)\subset\cM_{-1}$ for the one
flat class $\delta_{-1}$. The differential of $\Psi$ at $(X,0)$ is injective by Lemma~\ref{lem:detM}, and
$\dim\VZ+1=\dim\cM_{-1}$. So the image contains a neighbourhood of $X$ in $\cM_{-1}$; every $\theta_Q$ vanishes
on it, and $\NL(-1)=\cM_{-1}$.
\end{proof}

\begin{lemma}\label{lem:branch}
Let $(a,b)=(a_1+e_1,a_2+e_1)$ and let $V'$ be the locus of smooth surfaces of degree $d$ containing a complete
intersection of type $(a,b)$. The closure of $V'$ is irreducible of codimension $c_1+c_2-1$, and $\cM_{-1}$ is a
branch of it at $X$.
\end{lemma}
\begin{proof}
$V'$ is the image of an open subset of the incidence variety of pairs $(D,X'')$, a projective bundle over the
irreducible family of complete intersections of type $(a,b)$; so its closure is irreducible. By
Proposition~\ref{prop:minus1}, $\cM_{-1}=\Psi(U)$ and $\Psi(U\cap\{t\neq0\})\subset V'$; so $\cM_{-1}$ lies in the
closure and $\dim V'\ge\dim\cM_{-1}$. Conversely, at a pair $(D,X'')$ with $X''=\{n_1m_1+n_2m_2=0\}$ smooth and
$D=\{n_1=n_2=0\}$, the image of the differential of the projection to $S_d$ is $(n,m)_d$ (first lines of the
proof of Proposition~\ref{prop:first}(b); the kernel may be nonzero when $a+b\ge d$). Its codimension is
$\dim(S/(n,m))_d$, which depends only on $(a,b,d)$, because $(n_1,n_2,m_1,m_2)$ is a regular sequence (argument
of Lemma~\ref{lem:alg}(b) for $X''$). Moreover $(n,m)_d=T\,\NL([D])$ by Lemma~\ref{lem:pi} and
Proposition~\ref{prop:first}(a), whose proofs apply to $D\subset X''$ because the generators have degree $\le d$
and $[D]^2_{\mathrm{prim}}=-ab(a_1+e_2)(a_2+e_2)/d\ne0$. At $(D_t,X_t)$,
$t\ne0$, this codimension equals $\codim T\,\NL(-1)=c_1+c_2-1$, since $\NL([D_t])=\NL(-1)=\cM_{-1}$ near $X_t$. So the
projection has generic rank $\dim\cM_{-1}$, and $\dim V'=\dim\cM_{-1}$.
\end{proof}

\section{Consequences}\label{sec:cons}

\begin{proof}[Proof of Theorem A(5)]
Let $r\notin\{0,-1\}$. Since $dy(G_r)\ne0$, Proposition~\ref{prop:law} shows that $u'_Q(X)=0$ for all $Q\in K$ if
and only if $A=0$ at $X$; so the condition ``$A=0$'' does not depend on the choices ($\gt_i$, $y$, $Q'$) nor on
$r$. If $A\ne0$ at $X$, some $u'_Q$ is a unit and the ideal $(y^2u'_Q)_{Q\in K}$ equals $(y^2)$. In general
$\NL(r)_{\mathrm{red}}=\VZ\cup W_r$ with $W_r:=\{u'_Q=0\ \forall Q\in K\}$. The incidence variety $\cI^\circ$ of
triples $(Z_1,Z_2,X')$ ($X'$ smooth, $Z_1\cap Z_2=\emptyset$) is a connected complex manifold: an open subset of
a projective bundle over an open subset of $\cH_1\times\cH_2$. The set where $A=0$ is a closed analytic subset
(locally $\{u'=0\}$ on $\VZ$, by \S\ref{sec:struct} applied at nearby points). If it is not all of $\cI^\circ$ it
is nowhere dense; then near every $X$ the function $u'$ is not identically zero on $\VZ$, $W_r$ is a proper
analytic subset of $\cM_r$, and every component of $\NL(r)$ through $X$ has codimension $\ge c_1+c_2$.
\end{proof}

\begin{remark}\label{rem:rows}
(i) The second row of $N_t$ gives a complete intersection of type $(a_1+e_2,a_2+e_2)$ and class
$[Z_1]-[Z_2]+e_2H$; the columns give complete intersections of types $(a_1+e_1,a_1+e_2)$, $(a_2+e_1,a_2+e_2)$
and classes $[Z_2]-[Z_1]+a_1H$, $[Z_2]-[Z_1]+a_2H$ (Lemma~\ref{lem:Y} for $M^{\mathsf T}$). All have the same
Noether--Lefschetz locus. (ii) $\NL(-1)$ is a branch of the closure of the known family of surfaces containing a
complete-intersection curve (Lemma~\ref{lem:branch}). For $r\notin\{0,-1\}$, at a general point of $\VZ$ the only
component is $\VZ$ itself, in the types where $A\not\equiv0$ is established (Table~\ref{tab:types}). (iii) A
consistency check (referee~7; repeated by referee~8): Theorem~A(4) and (i) force $c(a,b;d)=c_1+c_2-1$ for the
four types $(a,b)$ of (i), with $c(a,b;d)$ the codimension of Proposition~\ref{prop:first}(a) for type $(a,b)$;
this holds in all $2352=588\times4$ cases with $d\le16$ ($588$ types, four complete intersections each; Hilbert
series).
\end{remark}

\section{The form \texorpdfstring{$A$}{A} on \texorpdfstring{$I(Z_1)I(Z_2)$}{I(Z1)I(Z2)}: Theorem A(6)}\label{sec:F}

Notation of \S\S\ref{sec:first}--\ref{sec:det}. Put $\rho_1:=\Res_S[\,\cdot\,/(f_1,f_2,g_1,g_2)]$ on $S_{2d-4}$,
so that $\pi_1=\kappa_1\rho_1$ (Lemma~\ref{lem:pi}) and $\rho_1(m\det M)=\varepsilon(m)$ for $m\in S_{d-4}$
(Lemma~\ref{lem:alg}(e)). In \eqref{eq:F}, $\adj(M)_{ji}=\pm M_{i'j'}$ has degree $a_i+e_j$, and the right-hand
side does not depend on the choice of the $m_{ij}$ or of $M$ (the changes alter $\operatorname{tr}(m\adj M)$ by
elements of $(f,h)$). Two facts are used throughout:
$\pi_1$ vanishes on $(f,g)_{2d-4}$ (Lemma~\ref{lem:pi}), and
\begin{equation}\label{eq:I2}
\varphi_1\text{ vanishes on }(I(Z_1)^2)_{3d-4},\qquad \varphi_2\text{ vanishes on }(I(Z_2)^2)_{3d-4}.
\end{equation}
Indeed, for $G,G'\in I(Z_1)_d$ the curve $Z_1$ lies in $\{F+tG+sG'=0\}$, so the $ts$-coefficient
$2\varphi_1(QGG')$ of $\Phi^{[Z_1]}_Q$ in \eqref{eq:exp} vanishes; and
$S_{d-4}\cdot I(Z_1)_d\cdot I(Z_1)_d=(I(Z_1)^2)_{3d-4}$.

\begin{lemma}[vector-field form of Griffiths--Dwork]\label{lem:vf}
Let $\xi=\sum_j\xi_j\partial_j$ be a vector field whose coefficients are forms of one degree $\nu$, let
$k\in\{1,2\}$, and let $m$ be a form of degree $2d-3-\nu-a_k$. Then
\begin{equation}\label{eq:ss}
\varphi_1\bigl(f_km\,(g_1\,\xi(f_1)+g_2\,\xi(f_2))\bigr)=\tfrac12\,\pi_1\bigl(\xi(f_k)\,m\bigr).
\end{equation}
Let $\Delta_{pq}=\partial_pf_1\,\partial_qf_2-\partial_qf_1\,\partial_pf_2$ $(0\le p<q\le3)$ and let
$J_2=(\Delta_{pq})$ be the ideal of the $2\times2$ minors of $\operatorname{Jac}(f_1,f_2)$. Then for
$k,l\in\{1,2\}$, $\mu\in(f)+J_2$ and every form $m$ with $\deg(\mu m)=2d-4-a_k+a_l$,
\begin{equation}\label{eq:sss}
\varphi_1(f_kg_l\,\mu m)=\tfrac12\,\delta_{kl}\,\pi_1(\mu m).
\end{equation}
\end{lemma}
\begin{proof}
Let $P=f_km$. Fact (G2) for the pole order $3$, applied to the forms $\xi_jP\in S_{2d-3}$ and summed over $j$,
gives $\xi(F)P\,\Omega/F^3\equiv\frac12(\xi(P)+\operatorname{div}(\xi)P)\,\Omega/F^2$; no further term appears.
Now $\xi(F)=\sum_i(g_i\,\xi(f_i)+f_i\,\xi(g_i))$, and $\varphi_1$ kills $f_kf_i\,m\,\xi(g_i)$ by \eqref{eq:I2}; on
the right, $\xi(P)+\operatorname{div}(\xi)P\equiv\xi(f_k)m$ modulo $(f)$, which $\pi_1$ kills. This is
\eqref{eq:ss}. The vector fields $\xi^1_{pq}=\partial_qf_2\,\partial_p-\partial_pf_2\,\partial_q$ and
$\xi^2_{pq}=-\partial_qf_1\,\partial_p+\partial_pf_1\,\partial_q$ satisfy
$\xi^l_{pq}(f_i)=\delta_{li}\Delta_{pq}$; so \eqref{eq:ss} for $\xi^l_{pq}$ is \eqref{eq:sss} for
$\mu=\Delta_{pq}$. For $\mu\in(f)$ both sides of \eqref{eq:sss} vanish, by \eqref{eq:I2} and because $\pi_1$
kills $(f)$.
\end{proof}

\run{Hypothesis (H) for $Z_1$.} \emph{$Z_1$ is smooth, and $S_j=((f)+J_2)_j$ for all $j\ge2d-4-|a_1-a_2|$.}
As $Z_1$ is smooth exactly when $S/((f)+J_2)$ is Artinian, (H) says that the top degree $t_0(Z_1)$ of this ring
satisfies $t_0(Z_1)<2d-4-|a_1-a_2|$. Hypothesis (H) for $Z_2$ is the same statement for $(h;e_1,e_2)$.

\begin{lemma}[reduction lemma, ``Lemma R'']\label{lem:R}
Assume \textup{(H)} for $Z_1$. For $\{k,k'\}=\{1,2\}$ and $N\in S_{3d-4-a_k}$,
\[
\varphi_1(f_kN)=\tfrac{\kappa_1}2\,\hat\psi_k(N),
\]
where $\hat\psi_k$ is the functional on $S_{3d-4-a_k}$ that vanishes on $(f_1,f_2,g_k^2,g_{k'})$ and satisfies
$\hat\psi_k(n\,g_k)=\rho_1(n)$ for $n\in S_{2d-4}$.
\end{lemma}
\begin{proof}
$(f_1,f_2,g_k^2,g_{k'})$ is a regular sequence with socle degree $3d-4-a_k$, and by (R) for
$(f,g_k^2,g_{k'})=\operatorname{diag}(1,1,g_k,1)\,(f,g_k,g_{k'})$ its residue sends $n\,g_k$ to $\pm\rho_1(n)$. So
$\hat\psi_k$ exists and is unique, and $S_{3d-4-a_k}=(f,g_k^2,g_{k'})_{3d-4-a_k}+g_kS_{2d-4}$. Let
$\lambda_k(N):=\varphi_1(f_kN)$. (a) $\lambda_k$ kills $(f)$, by \eqref{eq:I2}. (b) $\lambda_k(g_{k'}n)=0$ for
$n\in S_{2d-4-a_k+a_{k'}}$: $n\in(f)+J_2$ by (H), and \eqref{eq:sss} with $l=k'$. (c) $\lambda_k(g_k^2n)=0$ for
$n\in S_{d-4+a_k}$: $g_kn$ has degree $2d-4$, so $g_kn\in(f)+J_2$ by (H), and \eqref{eq:sss} with $l=k$ gives
$\lambda_k(g_k^2n)=\frac12\pi_1(g_kn)=0$. (d) $\lambda_k(g_kn)=\frac12\pi_1(n)=\frac{\kappa_1}2\rho_1(n)$ for
$n\in S_{2d-4}$, by (H) and \eqref{eq:sss}. By (a)--(d), $\lambda_k=\frac{\kappa_1}2\hat\psi_k$.
\end{proof}
For a line, $J_2=S$, (H) is empty, and Lemma~\ref{lem:R} is Lemma~\ref{lem:oct}(c) below, in any degree.

\begin{proof}[Proof of Theorem A(6)]
\emph{Step 1: \textup{(H)} holds for $Z_1$.} Let $Q=f_ih_jm$ and $\{i,i'\}=\{1,2\}$. The cross term of $A$
vanishes, because $Q\gt_1\gt_2\in I(Z_2)^2$ and \eqref{eq:I2}; so $A(Q)=\varphi_1(Q\gt_1^2)$. For $n\in S_{d-4}$,
$\pi_1(n(\kappa_1\gt_1-\det M))=\kappa_1\varepsilon(n)-\kappa_1\varepsilon(n)=0$, so by
Proposition~\ref{prop:first}(a) $\gt_1\equiv\kappa_1^{-1}\det M+s_1g_1+s_2g_2$ modulo $(f)$, and modulo
$(f_1,f_2,g_i^2,g_{i'})$ one has $\gt_1^2\equiv\kappa_1^{-2}(\det M)^2+2\kappa_1^{-1}\det M\cdot s_ig_i$. By
Lemma~\ref{lem:R}, $A(Q)=\frac{\kappa_1}2\hat\psi_i(h_jm\,\gt_1^2)$. The second term contributes
$\hat\psi_i(h_jm\det M\,s_ig_i)=\rho_1(h_jms_i\det M)=\varepsilon(h_jms_i)=0$. For the first,
$\adj(M)\,g=\det M\cdot h$, so $h_jm(\det M)^2=\sum_{l}m\det M\,\adj(M)_{jl}\,g_{l}$; $\hat\psi_i$ kills the
multiples of $g_{i'}$ and $\hat\psi_i(ng_i)=\rho_1(n)$, so
$\hat\psi_i(h_jm(\det M)^2)=\rho_1(m\,\adj(M)_{ji}\det M)=\varepsilon(m\,\adj(M)_{ji})$. Hence
$A(f_ih_jm)=\frac{\kappa_1}2\,\kappa_1^{-2}\,\varepsilon(m\,\adj(M)_{ji})$, which is \eqref{eq:F}.

\emph{Step 2: \textup{(H)} holds for $Z_2$.} The same computation with the roles exchanged
($F=h^{\mathsf T}M^{\mathsf T}f$) gives $B(f_ih_jm)=\frac1{2\kappa_2}\,\varepsilon'(m\,\adj(M^{\mathsf T})_{ij})$
with $\varepsilon'=\Res_S[\,\cdot\,/(h_1,h_2,f_1,f_2)]=\varepsilon$ (an even permutation) and
$\adj(M^{\mathsf T})_{ij}=\adj(M)_{ji}$. As $\kappa_1=\kappa_2$ (Lemma~\ref{lem:detM}) and $A=B$
(Proposition~\ref{prop:law}), this is \eqref{eq:F} again.

\emph{Step 3: continuity.} Let $\mathcal U$ be the set of $(f,h,M)$ with $Z_1\cap Z_2=\emptyset$ and $X$ smooth,
a non-empty Zariski-open subset of an affine space, hence connected. For fixed $m$ both sides of \eqref{eq:F}
are holomorphic on $\mathcal U$ ($\rho_1$ and $\varepsilon$ are regular; $\pi_1$ and the $\varphi_i$ are
single-valued periods, the cycles being given by $f$ and $h$; $\kappa_1=\pi_1/\rho_1$ is nowhere zero; $A$ does
not depend on the choice of the $\gt_i$), and (H) is a Zariski-open condition on $f$, resp.\ on $h$; the set of $(f,h,M)\in\mathcal U$ with
(H) for $f$ is non-empty, being the trace on $\mathcal U$ of a non-empty open subset of the irreducible affine
space. So
\eqref{eq:F} holds on all of $\mathcal U$ as soon as (H) holds for one general curve.

\emph{The last sentence of Theorem~\textup{A(6)}} (``only if''). If $e_j\ge4$, then $k_jS_{e_j-4}\subset K$ and
$k_j\notin I(Z_2)$, since otherwise $X$ is singular on $Z_2\cap\{k_{j'}=0\}$; symmetrically for $a_i\ge4$.
\end{proof}

\run{Where \textup{(H)} is proved, and where it is only computed.} By hand: (a) for a line, $J_2=S$; (b) for
a smooth conic, $(f)+J_2=(x_0,\dots,x_3)$ and $t_0=0$; (c) for a smooth plane curve of degree $a$,
$S/((f)+J_2)$ is the Milnor algebra of the curve and $t_0=3a-6$, so (H) holds if and only if
$3a-6<2d-4-(a-1)$; (d) every smooth curve of type $(2,2)$ is a diagonal pencil, $J_2=(x_px_q)_{p<q}$, and
$S/((f)+J_2)$ has Hilbert function $1,4,2$, so $t_0=2$ (case (d) was supplied by referee~7b). In each of the
$15$ types of Table~\ref{tab:types} with $K\neq0$ one of the two curves, taken general, is covered: seven types
contain a line, five a conic, $(1,3)+(1,3)$ and $(1,3)+(2,2)$ a plane cubic ($3<2\cdot8-4-2$), and
$(2,2)+(2,2)$ a $(2,2)$ curve ($2<12$). Only computed (\path{octic_main/t0_jacobian.py}, random curves modulo
$p$; reproduced by referee~7b's program): $t_0=5,8,7$ for the types $(2,3)$, $(2,4)$, $(3,3)$, and $12$ for
$(4,4)$ in referee~7b's program. These
values and (b)--(d) fit $t_0=2\min(a_1,a_2)+3\max(a_1,a_2)-8$; if this held for general curves of every type
other than $(1,1)$, then (H) would hold for one of the two curves in every type. But it is an observation modulo
$p$, not a theorem: for types without a line, a conic, a plane curve in the range of (c) or a $(2,2)$ curve,
\eqref{eq:F} is not proved.

\begin{corollary}\label{cor:F}
Assume the hypothesis of Theorem~\textup{A(6)}.
\begin{enumerate}[label=\textup{(\alph*)}]
\item If $\max(a_1,a_2)+\max(e_1,e_2)\ge4$, i.e.\ if $(I(Z_1)I(Z_2))_{d-4}\neq0$, then $A\not\equiv0$ on
$\cI^\circ$.
\item $A$ vanishes on $(I(Z_1)I(Z_2))_{d-4}$ at $X$ if and only if $M_{i'j'}\in(f,h)$ for every $(i',j')$ with
$a_{i'}+e_{j'}\ge4$. This locus can be empty among smooth surfaces: in type $(1,1)+(1,3)$ it forces
$g_1,g_2\in(f,h_1)$, and then $X$ is singular; in that type $K=(I(Z_1)I(Z_2))_2$, so $A\neq0$ at every point.
\item If $Z_1$ is the line $\{x_0=x_1=0\}$, then \eqref{eq:F} reads
$A(x_{i-1}h_jm)=\frac1{2\kappa_1}\,\varepsilon_L\bigl(m|_L\,\adj(M|_L)_{ji}\bigr)$ with
$\varepsilon_L=\Res_{R_L}[\,\cdot\,/(\bar h_1,\bar h_2)]$, where $R_L=S/I(L)$ and $\bar h_j=h_j|_L$; for
$\max(e_1,e_2)\le3$ this is all of $A$.
\end{enumerate}
\end{corollary}
\begin{proof}
(a) Choose $(i,j)$ with $a_i+e_j\le d-4$. For general $M$ the entry $\adj(M)_{ji}=\pm M_{i'j'}$ is an arbitrary
form of degree $a_i+e_j\le d-4$, not in $(f,h)$, and $\varepsilon(m\,\adj(M)_{ji})\neq0$ for a suitable $m$ by
the perfect pairing (R). A smooth $X$ with general $M$ exists: for smooth disjoint $Z_1,Z_2$ the map
$M\mapsto f^{\mathsf T}Mh$ is onto $I(Z_1\cup Z_2)_d$, and this ideal is generated by the $f_ih_j$, of degree
$\le d-2$; so the linear system $I(Z_1\cup Z_2)_d$ has base scheme $Z_1\cup Z_2$ and its general member is smooth
off $Z_1\cup Z_2$ (Bertini). At a point $p$ of $Z_i$ the differentials at $p$ of the members span the conormal
plane of $Z_i$ at $p$ (the ideal sheaf twisted by $d$ is globally generated), so being singular at $p$ is two
conditions; as $\dim(Z_1\cup Z_2)=1$, the general member is smooth. The conditions ``$X$ smooth'' and
``$M_{i'j'}\notin(f,h)$'' are both open and non-empty on the irreducible space of $M$. (b) By \eqref{eq:F} and (R), $A$ vanishes on
$(I(Z_1)I(Z_2))_{d-4}$ if and only if $\adj(M)_{ji}\in(f,h)$ whenever $a_i+e_j\le d-4$; the complementary
indices $(i',j')$ are those with $a_{i'}+e_{j'}\ge4$. In type $(1,1)+(1,3)$ the condition is
$M_{12},M_{22}\in(f,h_1)$; then $g_1,g_2\in(f,h_1)$ and $X$ is singular at the point $L\cap\{h_1=0\}$.
(c) For a line, $\varepsilon(N)=\varepsilon_L(N|_L)$.
\end{proof}

\run{Status of this section.} Working notes (\path{octic/PROOF_O_addendum.md}, sections C, D), read by one
AI referee in a follow-up: referee~7b (``no fatal error''; six small repairs, incorporated), and referee~6b
for the case of a line. Their own programs confirmed \eqref{eq:F} modulo $p$ on all generators $f_ih_jm$
(referee~7b: six types, in five of them with the constant $1/2\kappa_1$ and the sign, negative controls
failing; referee~6b: five types with a line). See \S\ref{sec:open}, item~(3).

\section{The octic case: a line and a disjoint \texorpdfstring{$(3,3)$}{(3,3)} curve}\label{sec:octic}

Here $d=8$, $Z_1=L=\{x_0=x_1=0\}$, $f_i=x_{i-1}$, $Z_2=C=\{h_1=h_2=0\}$ with cubics $h_j$; $M$ is a matrix of
quartics, $g=Mh$ (septics), $k=M^{\mathsf T}f$ (quintics); $T=S_8$ has dimension $165$ and $p_g=35$. Let
$R_L=S/I(L)=\CC[\sigma,\tau]$, and $\bar h_j$, $\Lambda$, $q_i$ as in Theorem~B. Since
$\nabla F|_L=(q_1,q_2,0,0)$, the septics $q_1,q_2$ are coprime; $\bar h_1,\bar h_2$ are coprime, and
$(\bar h)_j=(R_L)_j$ for $j\ge5$. Put $\varepsilon_L:=\Res_{R_L}[\,\cdot\,/(\bar h_1,\bar h_2)]$ on $(R_L)_4$,
$\rho_L:=\Res_{R_L}[\,\cdot\,/(q_1,q_2)]$ on $(R_L)_{12}$, $\rho_C:=\Res_S[\,\cdot\,/(h_1,h_2,k_1,k_2)]$ on
$S_{12}$, and on $(R_L)_{19}$
\[
\psi_1:=\Res_{R_L}[\,\cdot\,/(q_1^2,q_2)],\qquad\psi_2:=\Res_{R_L}[\,\cdot\,/(q_1,q_2^2)]
\]
(ordered denominators). Equivalently, $\psi_i$ vanishes on $(q_i^2,q_{i'})_{19}$ and $\psi_i(mq_i)=\rho_L(m)$,
where $\{i,i'\}=\{1,2\}$. Thus $\cD_{ij}=\psi_i(\bar h_j(\det\Lambda)^2)$.

\run{First order: Theorem B(1).} In degree $4$, $(f,g)_4=I(L)_4$ and $(h,k)_4=(h)_4$, so by
Propositions~\ref{prop:first} and \ref{prop:excess}: $c_1=\dim(R_L)_4=5$, $c_2=35-8=27$, and
$K=I(L)_4\cap(h)_4=\langle f_ih_j\rangle$ (if $l_1h_1+l_2h_2$ vanishes on $L$ then
$l_1|_L\bar h_1+l_2|_L\bar h_2=0$, so $l_j|_L=0$). For $f=(x_0,x_1)$ and forms $\varphi_3,\varphi_4$,
$\Res_S[N/(x_0,x_1,\varphi_3,\varphi_4)]=\Res_{R_L}[N|_L/(\varphi_3|_L,\varphi_4|_L)]$ (both are normalised by
the Jacobian). Hence $\varepsilon=\varepsilon_L(\,\cdot\,|_L)$, $\rho_1=\rho_L(\,\cdot\,|_L)$, $\kappa=\kappa_1$,
and $\kappa'=\kappa$ by Lemma~\ref{lem:detM} ($\kappa,\kappa'$ as in Lemma~\ref{lem:oct}(a) below); so
$A=\cD/2\kappa$ is \eqref{eq:F} with the normalisation of Lemma~\ref{lem:AB}. The incidence variety $\cI$ is a $\PP^{92}$-bundle over an
open subset of $G(1,3)\times G(2,S_3)$, irreducible of dimension $132$ (surfaces being counted in $\PP(S_8)$, of
dimension $164$; codimensions are the same in $S_8$ and in $\PP(S_8)$).

\begin{lemma}\label{lem:oct}
\begin{enumerate}[label=\textup{(\alph*)}]
\item $\pi_L=\kappa\,\rho_L(\,\cdot\,|_L)$ and $\pi_C=\kappa'\rho_C$ with $\kappa,\kappa'\ne0$.
\item $\varphi_L$ vanishes on $(I(L)^2)_{20}$ and $\varphi_C$ vanishes on $(I(C)^2)_{20}$.
\item $\varphi_L(f_i\tilde h)=\frac\kappa2\,\psi_i(\tilde h|_L)$ for $\tilde h\in S_{19}$, $i=1,2$.
\end{enumerate}
\end{lemma}
\begin{proof}
(a) The statement for $C$ is Lemma~\ref{lem:pi}. As there, $\pi_L$ kills $I(L)_{12}=S_4\cdot I(L)_8$. By (G2),
$\pi_L(m\,\partial_jF)=\int_L\Res(\partial_jm\,\Omega/F)=0$, and $\partial F/\partial f_i\equiv q_i$ mod $I(L)$.
So $\pi_L$ factors through the socle $(R_L)_{12}/(q)_{12}$; it is nonzero because
$\langle L,L\rangle_{\mathrm{prim}}=-6-\frac18$.
(b) For $G,G'\in I(L)_8$ the line lies in $\{F+tG+sG'=0\}$; the $ts$-coefficient of $\Phi_Q$ is
$2\varphi_L(QGG')=0$, and $S_4\cdot I(L)_8\cdot I(L)_8=(I(L)^2)_{20}$. The same argument applies to $C$.
(c) By (G2) with $l=2$, $\varphi_L(m\,\partial_jF)=\frac12\pi_L(\partial_jm)$. For $m=f_im'$ with
$m'\in(R_L)_{12}$: $f_im'\,\partial F/\partial f_j\equiv f_im'q_j$ mod $I(L)^2$, so
$\varphi_L(f_im'q_j)=\frac\kappa2\delta_{ij}\rho_L(m')$; and $\psi_i(m'q_j)=\delta_{ij}\rho_L(m')$ by (R). Since
$(R_L)_{19}=q_1(R_L)_{12}+q_2(R_L)_{12}$, the two functionals agree.
\end{proof}

\begin{proof}[Proof of Theorem B(2)]
Choose $\gt_1\in I(C)_8$ with $\gt_1|_L\equiv\kappa^{-1}\det\Lambda$ mod $(q)_8$ (possible since
$I(C)_8|_L=(R_L)_8$), and $\gt_2\in I(L)_8$ with $\kappa'\rho_C(Q\gt_2)=\varepsilon_L(Q|_L)$ for all $Q\in S_4$
(possible as in Lemma~\ref{lem:AB}). By (R) with $q=\Lambda\bar h$, $\rho_L(m\det\Lambda)=\varepsilon_L(m)$ for
$m\in(R_L)_4$, so $\pi_L(Q\gt_1)=\varepsilon_L(Q|_L)=\pi_C(Q\gt_2)$, while $\pi_L(Q\gt_2)=0=\pi_C(Q\gt_1)$. So
$G_r:=\gt_1-\gt_2/r\in T\,\NL(r)\setminus T\VZ$: this is the normalised excess vector, and we are in the
situation of Lemma~\ref{lem:AB}. For $Q\in K\subset I(L)\cap I(C)$ the forms $Q\gt_1\gt_2$, $Q\gt_2^2$ lie in
$I(L)^2$ and $Q\gt_1^2$, $Q\gt_1\gt_2$ lie in $I(C)^2$; by Lemma~\ref{lem:oct}(b),
\[
\varphi_\delta(QG_r^2)=\alpha(Q)+\beta(Q)/r,\qquad\alpha(Q)=\varphi_L(Q\gt_1^2),\quad\beta(Q)=\varphi_C(Q\gt_2^2).
\]
For $Q=f_ih_j$: $Q\gt_1^2\equiv f_i\cdot\bar h_j\gamma^2$ mod $I(L)^2$ with
$\gamma=\gt_1|_L=\kappa^{-1}\det\Lambda+s_1q_1+s_2q_2$. Modulo $(q_i^2,q_{i'})$,
$\gamma^2\equiv\kappa^{-2}(\det\Lambda)^2+2\kappa^{-1}\det\Lambda\cdot s_iq_i$, and
$\psi_i(\bar h_j\det\Lambda\,s_iq_i)=\rho_L(\bar h_js_i\det\Lambda)=\varepsilon_L(\bar h_js_i)=0$. So by
Lemma~\ref{lem:oct}(c), $\alpha(f_ih_j)=\frac\kappa2\kappa^{-2}\psi_i(\bar h_j(\det\Lambda)^2)=\cD_{ij}/2\kappa$.
Finally $\beta=\alpha$ by Proposition~\ref{prop:law}; the constant $\kappa'$ is never needed (in fact
$\kappa'=\kappa$ by Lemma~\ref{lem:detM}, for compatibly normalised residues).
\end{proof}

Theorem~B(3) is Theorem~A(4) with Lemma~\ref{lem:branch} for this type ($e_1=3$); here the ideal of $Y$ is
$I(L)\cap(h_1,k_2)$, because $k_2\in I(L)$, $I(L)$ is prime and $h_1\notin I(L)$, so $Y=L\cup C'$ as schemes. The
residue formula uses $K\subset I(L)\cap I(C)$. Among the types $(1,1)+(e,e)$, $d=2e+2$, this holds exactly for
$e\le3$ (and $K=0$ for $e=1,2$); for $e\ge4$ the forms $k_j$ enter $K$, and \eqref{eq:F} gives $A$ only on
$(I(L)I(C))_{d-4}$, although Theorem~A applies.

\subsection{The invariant is linear in \texorpdfstring{$M$}{M}; every smooth octic: Theorem B(4)}\label{ss:every}

\begin{lemma}[``Lemma L'']\label{lem:L}
$\cD_{ij}=\varepsilon_L\bigl(\adj(\Lambda)_{ji}\bigr)$ for $i,j\in\{1,2\}$; that is, since
$\adj\Lambda=\bigl(\begin{smallmatrix}\Lambda_{22}&-\Lambda_{12}\\-\Lambda_{21}&\Lambda_{11}\end{smallmatrix}\bigr)$,
\[
\cD=\begin{pmatrix}\cD_{11}&\cD_{12}\\\cD_{21}&\cD_{22}\end{pmatrix}
=\begin{pmatrix}\varepsilon_L(\Lambda_{22})&-\varepsilon_L(\Lambda_{21})\\
-\varepsilon_L(\Lambda_{12})&\varepsilon_L(\Lambda_{11})\end{pmatrix}.
\]
In particular $\cD$ is linear in $M$.
\end{lemma}
\begin{proof}
$\adj(\Lambda)\,q=\adj(\Lambda)\Lambda\bar h=\det\Lambda\cdot\bar h$, so
$\bar h_j(\det\Lambda)^2=\sum_l\det\Lambda\cdot\adj(\Lambda)_{jl}\,q_l$. The functional $\psi_i$ kills every
multiple of $q_l$, $l\neq i$, and $\psi_i(mq_i)=\rho_L(m)$ for $m\in(R_L)_{12}$. Hence
$\cD_{ij}=\rho_L(\det\Lambda\cdot\adj(\Lambda)_{ji})$, and $\rho_L(m\det\Lambda)=\varepsilon_L(m)$ for the binary
quartic $m=\adj(\Lambda)_{ji}$, by (R) for $q=\Lambda\bar h$.
\end{proof}
The working notes state this lemma with indices in $\{0,1\}$; the transcription above was made for
version~2; $\cD(X_*)$ is symmetric and cannot detect a transposition; the transcription was checked exactly at
random integer data with four distinct entries (referee~9: two data; the transposed formula and the formula with
the two residue symbols exchanged both fail there). The proof of Theorem~B(4) printed below replaces the
dimension count of the working notes. Exact checks:
six random integer data (\path{octic/check_D_linear.py}); referee~6b's own residue program at $11$ data.

\begin{corollary}\label{cor:L}
\textup{(a)} $\cD(X)=0$ if and only if all four entries of $\Lambda=M|_L$ lie in $(\bar h_1,\bar h_2)_4$, i.e.\ if
and only if $M\equiv0$ modulo $(x_0,x_1,h_1,h_2)$; this does not depend on the choice of $M$.

\textup{(b)} $\{\cD=0\}\subset\cI^\circ$ is irreducible \textup{(}and smooth\textup{)} of codimension exactly
$4$; it contains the split octics of Theorem~\textup{B(5)}.
\end{corollary}
\begin{proof}
(a) $\ker\varepsilon_L=(\bar h)_4$ by (R). The admissible changes of $M$ (after Lemma~\ref{lem:alg}) alter
$\Lambda$ by $w|_L\cdot(\bar h_2,-\bar h_1)$, with entries in $(\bar h)$. (b) Over the open set of disjoint pairs
$(L,C)$, with local frames $f,h$, let $\mathcal E$ be the vector bundle with fibre the $2\times2$ matrices of
quartics, and $\mu\colon\mathcal E\to\mathcal F$, $M\mapsto f^{\mathsf T}Mh$, the map onto the bundle
$\mathcal F$ with fibre $I(L\cup C)_8$ (onto by Lemma~\ref{lem:alg}(a)). The four linear forms
$M\mapsto\varepsilon_L(\Lambda_{ij})$ are independent ($\Lambda_{ij}=M_{ij}|_L$ is an arbitrary binary quartic and
$\varepsilon_L\ne0$) and vanish on $\ker\mu$ by (a); so they descend to four independent linear forms on the
fibres of $\mathcal F$. Their common zero set is intrinsic (transformation law below), hence a subbundle of
corank $4$. Now $\cI^\circ$ is the open subset of smooth surfaces in $\PP(\mathcal F)$, and $\{\cD=0\}$ is the
trace on it of the projectivised subbundle, which is irreducible and smooth; it is non-empty by
\S\ref{ss:split}.
\end{proof}

The form $\cD\in K^*$ is defined up to a nonzero scalar. It does not depend on $M$ (Corollary~\ref{cor:L}(a);
directly: $\det\Lambda$ changes by $q_2b_1-q_1b_2$, and
$\psi_i(\bar h_j\det\Lambda\,q_ib)=\varepsilon_L(\bar h_jb)=0$). Under $f\to Rf$, $h\to Ph$ one has
$\cD'=\lambda R\,\cD P^{\mathsf T}$ with $\lambda=1/(\det R\,(\det P)^2)$: indeed
$M'=R^{-\mathsf T}MP^{-1}$, so $\adj\Lambda'=P\adj(\Lambda)R^{\mathsf T}/(\det P\det R)$, and
$\varepsilon_L'=\varepsilon_L/\det P$ by (R). Locally
on $\cI^\circ$ the data $f,h,M$ are regular, so $\{\cD=0\}$ is closed. By Lemma~\ref{lem:struct} at the points of
$\VZ$ and Theorem~B(2): $u'=0$ at $X'\in\VZ$ if and only if $\cD(X')=0$.

\begin{proof}[Proof of Theorem B(4)]
Let $r\notin\{0,-1\}$. If $\cD(X)\neq0$, then $A\ne0$ at $X$ and $\NL(r)=\{y^2=0\}$ by Theorem~A(5). Let
$\cD(X)=0$. On the linear family $\Sigma=F+I(L\cup C)_8\subset\VZ$ (same curves, $K$ constant) the matrix $M$
varies by an arbitrary matrix $M'$ of quartics, so $\delta\Lambda=M'|_L$ is arbitrary, and by
Lemma~\ref{lem:L} $\cD$ is affine-linear on $\Sigma$, with linear part
\[
d\cD(\delta\Lambda)=\operatorname{cof}\bigl(\varepsilon_L(\delta\Lambda)\bigr),\qquad
\operatorname{cof}\begin{pmatrix}\mu_{11}&\mu_{12}\\\mu_{21}&\mu_{22}\end{pmatrix}
=\begin{pmatrix}\mu_{22}&-\mu_{21}\\-\mu_{12}&\mu_{11}\end{pmatrix},
\]
of rank $4$ at every point, because $\delta\Lambda\mapsto\varepsilon_L(\delta\Lambda)$ is onto the $2\times2$
matrices. By Lemma~\ref{lem:struct} and Theorem~B(2), $u'_Q|_\Sigma=E\cdot\cD(Q)$ with
$E(X')=(r+1)/(2r\kappa_{X'}\,dy_{X'}(G_{r,X'})^2)$, a nowhere-zero holomorphic function ($\kappa_{X'}$ is the
ratio of the nonzero functionals $\pi_L$ and $\rho_L$, which vary holomorphically; $G_{r,X'}$ solves a linear
system of constant rank; $y$ is fixed; $E$ does not depend on the choice of $G_{r,X'}$ modulo $T\VZ$, and only
its continuity at $X$ and its non-vanishing are used). As $\cD(X)=0$, $du'(X)=E(X)\,d\cD$ on $T\Sigma$, of rank $4$ on
$T\Sigma\subset T\VZ$. So the four functions $u'_Q$ have independent differentials on $\cM_r$, and so have
their restrictions to $\VZ$: $W_r$ is smooth of codimension $4$ in $\cM_r$, $W_r\cap\VZ$ is smooth of
codimension $4$ in $\VZ$ and equals $\{\cD=0\}$ near $X$, the ideal of $\NL(r)$ is $(y^2)\cdot(u')=(y^2)\cap(u')$
(the ideal $(u')$ is prime and does not contain $y$), and $\theta_Q=y^2u'_Q$ has leading term $dy^2\cdot du'_Q$,
of order $3$. Finally, at a point of $W_r\cap W_{r'}$, $r\ne r'$, both
$\delta_r$ and $\delta_{r'}$ are of type $(1,1)$, hence so are $[L]$ and $[C]$; as $\NL([L])=V_1$ and
$\NL([C])=V_2$ (Proposition~\ref{prop:first}), the point lies in $\VZ$.
\end{proof}

\emph{An exact point} (no longer needed for Theorem~B(4); kept as a check). Let $X_*=\{f^{\mathsf T}Mh=0\}$
with $h=(x_2^3+x_0^3,\ x_3^3+x_1^3)$ and
$M=\bigl(\begin{smallmatrix}x_0^4+x_2^4& x_2^2x_3^2+2x_2x_3^3\\ x_2^2x_3^2-x_2^3x_3& x_1^4+x_3^4\end{smallmatrix}\bigr)$;
it is smooth (a rank computation modulo several primes, also by referees~6 and~8), and
$\bar h=(\sigma^3,\tau^3)$. With the normalisation $\Res[\operatorname{Jac}(q_1,q_2)/(q_1,q_2)]=49$ the residues
of Theorem~B(2) give, exactly over $\QQ$, $\cD(X_*)=\bigl(\begin{smallmatrix}0&-1\\-1&0\end{smallmatrix}\bigr)$
(reproduced by referees~6 and~8 with independent programs). Lemma~\ref{lem:L} gives the same: $\varepsilon_L$ is
the coefficient of $\sigma^2\tau^2$, which is $1$ in $\Lambda_{12}$ and $\Lambda_{21}$ and $0$ in
$\Lambda_{11}=\sigma^4$ and $\Lambda_{22}=\tau^4$.

\subsection{Split octics and the Fermat point: Theorem B(5)}\label{ss:split}
Let $F=P(x_0,x_1)+\tilde P(x_2,x_3)$ be as in Theorem~B(5). Here the line is $L=\{\phi_1=\gamma_1=0\}$;
Theorem~B(1)--(4) apply with $f=(\phi_1,\gamma_1)$ in place of $(x_0,x_1)$, and $\sigma,\tau$ are the
restrictions to $L$ of a linear form in $x_0,x_1$ and of one in $x_2,x_3$ that do not vanish on $L$ (we write
$x_1,x_3$). One has
$F=(\phi_1,\gamma_1)\operatorname{diag}(\phi_4,\gamma_4)(\phi_3,\gamma_3)^{\mathsf T}$, and on $L$ every $\phi_j$
restricts to a multiple of $\sigma^j$ and every $\gamma_j$ to a multiple of $\tau^j$. After rescaling the
$h_j$: $\bar h=(\sigma^3,\tau^3)$ and $\Lambda_0=\operatorname{diag}(a\sigma^4,b\tau^4)$ with $ab\ne0$. Both
entries lie in $(\sigma^3,\tau^3)$, so $\cD=0$ by Corollary~\ref{cor:L}(a), and Theorem~B(4) gives
$\NL(r)=2\VZ\cup W_r$. Here $\varepsilon_L$ is the coefficient of $\sigma^2\tau^2$, so
$d\cD(\delta\Lambda)=\operatorname{cof}\bigl([\sigma^2\tau^2]\,\delta\Lambda\bigr)$; a direction with
$du'\ne0$ is $M'_{11}=x_1^2x_3^2$, i.e.\ $G=\phi_1\phi_3\,x_1^2x_3^2$. The data of \cite{Mo26} are of split form, as read by referee~6:
$x_0^8+x_1^8=\phi_1\phi_3\phi_4$, $x_2^8+x_3^8=\gamma_1\gamma_3\gamma_4$.

Modulo $p$ only (\path{octic/fermat3.py}; two primes, one program, examined by no referee), a direct
third-order computation at the Fermat octic agrees: the form of order $3$ on $T\,\NL(r)$ is
$dy^2\cdot\frac{r+1}{r}\hat\ell$ with $\operatorname{rank}\hat\ell=4$ on $T\VZ$. Theorem~B does not use it.

\begin{remark}[symmetry; found by referee 6, by hand]\label{rem:sym}
A split octic is invariant under $\operatorname{diag}(\mu,\mu,\nu,\nu)$, $\mu^8=\nu^8=1$, which fixes $L$ and
$C$. Let $\omega=\mu/\nu$. For $Q$ of bidegree $(j,4-j)$ in $(x_0,x_1\,|\,x_2,x_3)$ the form $Q\Omega/F$ has
weight $\omega^{j+2}$, so the function $\Phi_Q$ has weight $\omega^{-(j+2)}$; a deformation of bidegree
$(i,8-i)$ has weight $\omega^i$, and $y$ has weight $\omega^4$. The basis elements of $K$ have $j=4,1,3,0$.
Hence $u'_Q(X)=0$ (a second proof of $\cD=0$), and $du'_Q$ can be nonzero only in bidegrees $i\equiv6-j$ mod
$8$, as the formula for $d\cD$ confirms. Moreover, for the subgroup $\{\omega^4=1\}$ the weight of
$\theta_Q=y^2u'_Q$ is non-trivial for $Q\in K$, while $y\not\equiv0$ on the smooth fixed germ
$\cM_r^{\mathrm{fix}}$ (tangent to the bidegrees $(8,0),(4,4),(0,8)$): so $\cM_r^{\mathrm{fix}}\subset W_r$. A
slice $\mathcal S$ spanned by monomials and complementary to $T\VZ$ meets $\cM_r$ inside $\cM_r^{\mathrm{fix}}$:
indeed, $\mathcal S$ and $\cM_r$ (with bihomogeneous $Q'$) are invariant, so $\Gamma=\mathcal S\cap\cM_r$ is an
invariant smooth curve; its tangent line maps isomorphically and equivariantly onto $T\cM_r/T\VZ$, which has
weight $\omega^4$; hence the subgroup $\{\omega^4=1\}$ acts trivially on the tangent line and, being finite,
trivially on $\Gamma$. So on such a slice the locus is a smooth curve to all orders, for every $r$; the
obstruction lives in the bidegrees $(6,2),(5,3),(3,5),(2,6)$. This explains the slice results of
\cite[Thm~2]{Mo21} and \cite{Mo26} provided their slices are spanned by monomials. As read by machine, family
(9) of \cite{Mo21} is the Fermat octic minus a combination of $32$ listed monomials, and the slices of
\cite{Mo26} are spanned by monomials of bidegrees $(4,4)$ and $(8,0)$; a person should confirm this. If a slice
were not of this kind, this explanation would not apply. On any linear complement $\mathcal S$ of $T\VZ$ the
locus is $\{t^2u'_Q=0\}$ on the smooth curve $\Gamma=\mathcal S\cap\cM_r$ ($t=y|_\Gamma$): either $\Gamma$, or a
point of multiplicity $2+k$, $k=\min_Q\operatorname{ord}_0(u'_Q|_\Gamma)\ge1$. An $N$-smoothness statement on
$\mathcal S$ is contradicted by Theorem~B(5) exactly when $k\le N-2$ (for $N=5$: $k\le3$), for instance when
some $du'_Q$ does not vanish on $T\Gamma$; we have not computed $k$ for any slice.
\end{remark}

\subsection{Comparison with \texorpdfstring{\cite{Mo26}}{Movasati's preprint}}
Table~\ref{tab:mov} compares Theorem~B with the statements of \cite{Mo26}. Every statement in the left column is
his \emph{as we read it, through machine extraction}; the quotations were not compared with the PDF by a person.
His Conjecture~1 and the sentence of his abstract on ``all but finitely many $r$'' are conjectural statements;
his Propositions 2--5 and his tangent-space numbers are not contradicted, the slice statements on the condition
that his slices are spanned by monomials.

\begin{table}[!t]
\centering\footnotesize
\begin{tabular}{@{}p{0.47\textwidth}p{0.49\textwidth}@{}}
\toprule
Statement of \cite{Mo26}, as we read it (machine extraction) & Status here\\
\midrule
\rg Tangent codimensions $5,27,31,32$; $V_{C_1},V_{C_2}$ smooth and transversal, intersection of codimension $32$ &
 \rg confirmed: Theorem B(1), proved\tabularnewline
\rg Proposition 3: the locus is $2$-smooth in the full family & \rg confirmed: $\cD=0$ at the Fermat point,
 Theorem B(5)\tabularnewline
\rg Proposition 2 (conditional), on the $25$-dimensional family: infinitely many components as $r$ varies,
 without a codimension claim; he notes that they might be of codimension $35$ & \rg not contradicted: that
 family lies in the fixed space of Remark~\ref{rem:sym}, and the fixed germ lies in $W_r$; the components
 through the Fermat point other than $\VZ$ are the $W_r$, of codimension $35$\tabularnewline
\rg Proposition 4: on the $25$-dimensional family the loci form a ``general pencil'' (his term; as we read
 \cite{Mo21}: tangent codimension $p_g$ for some $r$, and no inclusion between the two tangent spaces).
 Proposition 5: $6$-smooth and $9$-smooth on small slices. Quoted from \cite[Thm~2]{Mo21}: $5$-smooth on the
 $32$-dimensional slice $\mathcal S$ & \rg not contradicted if $\mathcal S$ is spanned by monomials: the
 $25$-dimensional family and the small slices lie in the fixed space of the bidegrees $(8,0),(4,4),(0,8)$, and
 a monomial slice complementary to $T\VZ$ meets $\cM_r$ inside the fixed germ, which lies in $W_r$
 (Remark~\ref{rem:sym}); there the locus is smooth to all orders for every $r$\tabularnewline
\rg Introduction: a statement about a ``full theoretical proof of the smoothness'' of the locus on the
 $32$-dimensional slice $\mathcal S$ & \rg noticed by a machine reading (referee~9) only; it could not be
 interpreted (a claim, or a statement of what is missing); it must be read by a person, and no status is given
 here\tabularnewline
\rg The use of slice smoothness as evidence for smoothness in $T$ (an unbounded order appears only as the answer
 of a language model on a $10$-dimensional slice, which he partly disowns; the exact wording should be re-read by a
 person) & \rg the inference fails for $r\ne-1$: in $T$ the order of smoothness at the Fermat point is exactly
 $2$ (Theorem B(4), (5)). The slice statements themselves are not contradicted, on the same
 condition\tabularnewline
\rg Conjecture 1: for rational $r\neq0$ the local loci are smooth, distinct analytic spaces of codimension $31$
 through the common locus of codimension $32$ & \rg holds for $r=-1$ only (Theorem B(3)); for every other $r$,
 at every point of $\VZ$, the locus is not smooth and has no component of codimension $31$
 (Theorem B(4))\tabularnewline
\rg ``In particular'': the underlying varieties are distinct, of codimension $<35$ & \rg where $\cD\ne0$ the
 reduced loci all equal $\VZ$; where $\cD=0$, e.g.\ at the Fermat point, they are distinct, because
 $W_r\not\subset\VZ$ and the $W_r$ differ; components: $\VZ$ ($32$, the same for all $r$) and $W_r$
 ($35$)\tabularnewline
\rg Abstract: a possible counterexample to Harris's conjecture & \rg not supported by these classes at the
 points of $\VZ$: the only component of codimension $<35$ is $\VZ$ (and, for $r=-1$, the known locus of octics
 containing a $(4,4)$ complete intersection); the $W_r$ have the general codimension. Components that do not
 meet $\VZ$, and Harris's conjecture itself, are outside our statements\tabularnewline
\bottomrule
\end{tabular}
\caption{Theorem B and the statements of \cite{Mo26}. ``$N$-smooth'' is his terminology (\S\ref{sec:intro}).}\label{tab:mov}
\end{table}

\run{A proposed test for his code.} Add $G=\phi_1\phi_3\,x_1^2x_3^2$ ($\phi_1,\phi_3$ his factors of
$x_0^8+x_1^8$) to his $32$-dimensional slice $\mathcal S$, assumed to be a linear complement of $T\VZ$. On
$\mathcal S\oplus\CC G$ Theorem~B predicts a smooth surface for $r=-1$, and a locus that is $2$-smooth but not
$3$-smooth for $r\ne-1$. We have not run this test.

\section{Computations}\label{sec:comp}

Paths are relative to \path{night/} (Lean files: \path{lean/}; outputs: \path{octic_main/out/}; working
notes: \path{octic/PROOF_O2.md}, \path{PROOF_O.md}, \path{PROOF_O_addendum.md}; referees:
\path{referee*_report.md}, \path{referee6/}, \path{referee7/}, \path{referee8/}).

\begin{table}[!t]
\centering\footnotesize
\begin{tabular}{@{}clccccllc@{}}
\toprule
 & & & & & & \multicolumn{3}{c}{supports for $A\not\equiv0$, kept apart}\\
\cmidrule(l){7-9}
$d$ & types & $c_1,\ c_2$ & $\codim\VZ$ & $\codim T\,\NL(r)$ & $\dim K$ & proved; (H) for &
 exact; Lean & mod $p$\\
\midrule
4 & $(1,1)+(1,1)$ & $1,\ 1$ & 2 & 1 & 0 & \multicolumn{3}{l}{-- ($K=0$)}\\
5 & $(1,1)+(1,2)$ & $2,\ 3$ & 5 & 4 & 0 & \multicolumn{3}{l}{-- ($K=0$)}\\
6 & $(1,1)+(2,2)$ & $3,\ 8$ & 11 & 10 & 0 & \multicolumn{3}{l}{-- ($K=0$)}\\
6 & $(1,1)+(1,3)$ & $3,\ 6$ & 9 & 8 & 2 & the line & exact ($\Ob$) & yes\\
6 & $(1,2)+(1,2)$ & $5,\ 5$ & 10 & 9 & 1 & a conic & exact ($\Ob$); Lean & yes\\
7 & $(1,1)+(1,4)$ & $4,\ 9$ & 13 & 12 & $8^\dagger$ & the line & exact ($\Ob$) & yes\\
7 & $(1,1)+(2,3)$ & $4,\ 15$ & 19 & 18 & 2 & the line & exact ($\Ob$) & yes\\
7 & $(1,2)+(1,3)$ & $7,\ 9$ & 16 & 15 & 5 & the conic & exact ($\Ob$) & yes\\
7 & $(1,2)+(2,2)$ & $7,\ 12$ & 19 & 18 & 2 & the conic & exact ($\Ob$) & yes\\
8 & $(1,1)+(1,5)$ & $5,\ 12$ & 17 & 16 & $19^\dagger$ & the line & -- & yes\\
8 & $(1,1)+(2,4)$ & $5,\ 23$ & 28 & 27 & $8^\dagger$ & the line & -- & yes\\
8 & $(1,1)+(3,3)$ & $5,\ 27$ & 32 & 31 & 4 & the line & exact ($\cD$) & yes\\
8 & $(1,2)+(1,4)$ & $9,\ 13$ & 22 & 21 & $14^\dagger$ & the conic & -- & yes\\
8 & $(1,2)+(2,3)$ & $9,\ 21$ & 30 & 29 & 6 & the conic & -- & yes\\
8 & $(1,3)+(1,3)$ & $12,\ 12$ & 24 & 23 & 12 & a plane cubic & -- & yes\\
8 & $(1,3)+(2,2)$ & $12,\ 16$ & 28 & 27 & 8 & the plane cubic & -- & yes\\
8 & $(2,2)+(2,2)$ & $16,\ 16$ & 32 & 31 & 4 & a $(2,2)$ curve & -- & yes\\
10 & $(1,1)+(4,4)$ & $7,\ 62$ & 69 & 68 & $16^\dagger$ & the line & -- & yes\\
\bottomrule
\end{tabular}
\caption{All types with $d\le8$, and one with $d=10$. Numerical columns: Propositions \ref{prop:first},
\ref{prop:excess} (proved; recomputed by referees~7 and~8). The last three columns are separate supports for
$A\not\equiv0$: \emph{proved} by Theorem~A(6), with the curve for which (H) is proved by hand; \emph{exact} over
$\QQ$ at one explicit surface, resp.\ \emph{Lean}-certified finite computation at one sextic; \emph{modulo $p$}
at one random point ($\Ob_r\ne0$ for $r\in\{1,2,3,\frac12,-2,\frac53\}$, $\Ob_{-1}=0$). The entries ``$\Ob$'',
``Lean'' and ``mod $p$'' are conditional on the $\Ob$-link (``Lean'' also on smoothness and disjointness); ``mod $p$'' is not lifted to characteristic $0$.
$^\dagger$: $K$ is larger than $(I(Z_1)I(Z_2))_{d-4}$, on which alone \eqref{eq:F} describes $A$.}\label{tab:types}
\end{table}

\run{Modulo $p$} (\path{octic_main/second_order.py}; outputs \path{ci_*_s0.txt}, among them
\path{ci_1_1_3_3_s0.txt}). For each type: random curves and a random $M$ modulo $p=2147483629$, seed $0$; the
Jacobian-ring form $\Ob_r$ of the second-order invariant on $K$ (an expression which we take to be of the kind of the
second-order invariant attributed to Maclean (\cite{Mac}, not opened)), for $r\in\{1,-1,2,3,\frac12,-2,\frac53\}$. This is evidence at one random point per type
and one prime (further seeds or a second prime for $(1,2)+(1,2)$, $(1,1)+(1,3)$, $(2,2)+(2,2)$); to conclude in
characteristic $0$ one would still have to argue that the ranks at the point are the generic ranks, which is
not done. \emph{The law $(r+1)/r$} (outputs \path{law_ci_*_s3.txt}, \path{law_octic_random_s4.txt}): with $G_r$
normalised by a fixed class and a fixed basis of $K$, the vector $\Ob_r(G_r)\cdot\frac{r}{r+1}$ is the same for
all six tested $r\ne-1$, and $\Ob_{-1}=0$, in the types $(1,2)+(1,2)$, $(1,1)+(1,3)$, $(1,1)+(1,4)$,
$(1,2)+(2,2)$, $(1,1)+(3,3)$. Referee~7's own program ($p=2^{31}-1$, own derivation) found the same for seven
values of $r$ in the first three of these types; in $(1,1)+(1,4)$ the cross terms of Remark~\ref{rem:cross} do
not vanish separately, so this also tests $B=A$. \emph{Formula \eqref{eq:F}} (\texttt{FORMULA=1}; outputs
\path{formula_*_s7.txt}, \path{formula_octic_s11.txt}): in the $14$ types with $d\le8$ and $K\ne0$ the vector of
the values of the second-order form on all generators $f_ih_jm$ of $(I(Z_1)I(Z_2))_{d-4}$ is nonzero and
proportional to $(\varepsilon(m\,\adj(M)_{ji}))$.

\run{Exact over $\QQ$} (\path{octic_main/exact_sextic.py}; rational arithmetic with python-flint, no fitted
correction term; the explicit integer data are printed in the output files \path{exact_*.txt}). For the six
types $(1,1)+(1,3)$, $(1,2)+(1,2)$ ($d=6$) and $(1,1)+(1,4)$, $(1,1)+(2,3)$, $(1,2)+(1,3)$, $(1,2)+(2,2)$
($d=7$)---all types with $d\le7$ and $K\ne0$---each at one explicit surface with small integer coefficients
(seed $0$): $X$ is smooth, the curves are disjoint, $c_1$, $c_2$, $\dim K$ are as in Table~\ref{tab:types},
$\Ob_{-1}=0$, and $\Ob_r\cdot\frac{r}{r+1}$ is the same nonzero vector for $r=1,2,3,\frac12,-2,\frac53$.

\run{What ``exact'' means here: the $\Ob$-link.} It means exact evaluation of the Jacobian-ring expression
$\Ob_r(Q,G)=\Res_J\bigl(Q\,(G\operatorname{div}\xi-\xi\cdot\nabla G)\bigr)$, where $\Res_J$ is the socle
functional of the Jacobian ring, $P_rG=\xi\cdot\nabla F$, and
$P_r=\det\operatorname{Jac}(f_1,f_2,g_1,g_2)+r\det\operatorname{Jac}(h_1,h_2,k_1,k_2)$. That $\Ob_r$ is a constant
multiple of the form \eqref{eq:hess}---\emph{the $\Ob$-link}---rests on an argument which referee~7 re-derived
and tested with its own code in three types (with negative controls), but which is \emph{not written out as a
proof}: for an exact representative $P$ ($\Res(P\Omega/F^2)=\delta_{\mathrm{prim}}$) and $GP=\xi\cdot\nabla F$
one has $\varphi_\delta(QG^2)=c'\Res_J(Q(G\operatorname{div}\xi-\xi\cdot\nabla G))$ for $Q\in K$; and
$c\cdot\det\operatorname{Jac}(f,g)$ is an exact representative, because it is correct modulo $F^2H^2$ by the
period formula for complete intersections (\cite{CG,VL}; to be cited precisely),
$\Res(\det\operatorname{Jac}\cdot\Omega/F^2)$ is flat along the incidence variety, and a flat section of $F^2$
over $V_i$ vanishes since $(J,P)_{d-4}=0$. The exact runs, the Lean certificate and the runs modulo $p$ all
depend on this link; Theorem~A(6) and Theorem~B do not. The exact runs were made after referee~7's reading;
referee~8 read their outputs but did not re-run them.

\run{Lean-certified finite computation} (\path{lean/SexticCert.lean}: Lean~4 core only, no Mathlib, no
\texttt{native\_decide}, no \texttt{sorry}; $74$ theorems; re-run by the coordinating thread). At the sextic
$X=\{f^{\mathsf T}Mh=0\}$ of type $(1,2)+(1,2)$ with
\[
f=(x_0,\ x_2^2+x_1x_3),\qquad h=(x_1,\ x_3^2+x_0x_2),\qquad
M=\begin{pmatrix}x_1^4-x_0^4&-x_3^3-x_0x_1x_2\\-x_2^3&x_2x_3\end{pmatrix},
\]
and $Q_0=f_1h_1=x_0x_1$, the kernel checks the whole computation of $\Ob$: the definitions of $g$, $k$,
$\nabla F$ and of the two Jacobian determinants in $P_r$; explicit $\gt_1\in(h)$, $\gt_2\in(f)$ and vector fields
$B^0,B^1,B^2$; an explicit nonzero functional $\rho$ on $S_{16}$ vanishing on the Jacobian ideal and on
$Q_0\cdot\det\operatorname{Jac}\cdot S_6$ for both Jacobian determinants; and, for \emph{every} rational
$r\neq0$, with $G_r=\gt_1-\gt_2/r$ and $B_r=B^0-B^1/r+rB^2$, the identities $P_rG_r=B_r\cdot\nabla F$ and
$\Ob(r):=\rho\bigl(Q_0(B_r\cdot\nabla G_r-G_r\operatorname{div}B_r)\bigr)=A_0(1+\frac1r)$ with an explicit integer
$A_0\neq0$ (the opposite sign convention to that of $\Ob_r$ above; immaterial for the statements); so $\Ob(r)=0$ if and only if $r=-1$. \emph{Not} certified: smoothness of $X$, disjointness of the
conics, $\dim(S/J)_{16}=1$ (so that $\rho$ is the socle functional and $\dim K=1$; these are exact computations
in Python), the independence of the lift $B_r$, and the $\Ob$-link. This is a certified finite computation, not
a certified statement about Noether--Lefschetz loci. For comparison (not Lean-certified): at this sextic
\eqref{eq:F} gives $A(x_0x_1)=\varepsilon(x_2x_3)/2\kappa_1\neq0$, because $x_2x_3$ spans the socle of
$\CC[x_2,x_3]/(x_2^2,x_3^2)$; this agrees with $A_0\neq0$.

\run{Split pairs on Fermat surfaces} (\path{octic_main/split_fermat.py}, outputs \path{split_fermat_*.txt};
through $\Ob$, modulo one prime). On the Fermat surface of degree $d$ let
$Z_1=\{\phi_{a_1}(x_0,x_1)=\gamma_{a_2}(x_2,x_3)=0\}$ and $Z_2=\{\phi_{e_1}=\gamma_{e_2}=0\}$, the $\phi$, $\gamma$
being products of linear factors of $x_0^d+x_1^d$, resp.\ $x_2^d+x_3^d$; then
$M=\operatorname{diag}(\phi_{\mathrm{rest}},\gamma_{\mathrm{rest}})$. Call the arrangement balanced if $a_1+e_1=a_2+e_2=d/2$. With consecutive roots (first three rows):
\begin{center}\footnotesize
\begin{tabular}{@{}cll@{}}
\toprule
$d$ & arrangement $(a_1,a_2)+(e_1,e_2)$ & second-order form at the Fermat point\\
\midrule
6 & $(1,2)+(1,2)$;\ \ $(1,1)+(1,3)$, $(1,3)+(1,1)$ & nonzero for $r\ne-1$\\
8 & $(1,1)+(3,3)$;\ \ $(2,2)+(2,2)$;\ \ $(1,3)+(3,1)$\quad (balanced) & zero for all
 $r$\\
8 & $(1,3)+(1,3)$;\ \ $(1,2)+(2,3)$, $(1,1)+(2,4)$, $(1,2)+(1,4)$ & nonzero for $r\ne-1$\\
6 & $(1,2)+(1,2)$, the quadratic factors $\gamma$ being antipodal pairs of roots & zero for all $r$\\
\bottomrule
\end{tabular}
\end{center}
The first balanced case is Theorem~B(5). By \eqref{eq:F} (in the types where \eqref{eq:F} is proved; all
arrangements of the table are of such types): $M\equiv0$ modulo $(f,h)$ whenever
$|(a_1+e_1)-(a_2+e_2)|\le1$ (both entries then have degree above the socle degree); otherwise $A$ at a split
point is a product of binary residues, nonzero for general roots but zero for special ones. In the last row
$A=0$ both numerically and by \eqref{eq:F} (three quadratic forms in $x_2^2,x_3^2$ are dependent), while for
generic pairs of roots both are nonzero (\path{split_fermat_6_1212_antipodal.txt},
\path{split_fermat_6_1212_generic.txt}). So the vanishing at the Fermat point is not confined to the balanced
arrangements. \emph{Relation to \cite{Mo21}, as we read it} (machine digests): his Theorem~1 (singular at the Fermat point,
for coefficients in a stated range, in the cases other than the octic one with $4\le d\le8$) agrees with the
table in the unbalanced arrangements; in the balanced ones $(2,2)+(2,2)$, $(1,3)+(3,1)$ a singularity follows
from Theorem~A(5), (6): in the fifteen types, for $r\notin\{0,-1\}$, the ideal $(y^2u'_Q)_{Q\in K}$ with
$u'\not\equiv0$ is not radical, so $\NL(r)$ is non-reduced, hence singular, at every point of $\VZ$, the
balanced Fermat points included. His Remark~7 (``possibly smooth'' only for $r=0,\pm1$ in those cases): for
disjoint pairs with $d=a_1+a_2+e_1+e_2$ the locus is smooth for $r=-1$ by Theorem~A(4). His Remark~6: in all
seven exceptional cases of degree $6$ listed there the two curves share factors, i.e.\ are not disjoint, so
they are outside Theorem~A.

\run{Octic programs} (\path{octic/proofO_exact.py}, \path{check_D_linear.py}, \path{octic.py},
\path{fermat3.py}). Exact: $\cD(X_*)$; Lemma~\ref{lem:L} at six random integer data; the identities of
Lemma~\ref{lem:lean} over $\mathbb Z[t]$ at $X_*$; $\cD=0$ and $d\cD$ at the split point over $\QQ(a,b)$. Modulo
two primes: the codimensions $5,27,32,31$; the Jacobian-ring second-order matrix is proportional to $\cD$ for
$r=2,5$ and zero for $r=-1$. Referee~6's own code ($p=2^{31}-1$) found the Hessian proportional to $\cD$ with
factor $(r+1)/r$ for five values of $r$, and zero for $r=-1$, at $X_*$, at two dense random octics, at a split
and at a perturbed split octic; and, in its follow-up, $\cD=0$ and a zero second-order form at a non-split dense
octic with $M\equiv0$ modulo $(x_0,x_1,h_1,h_2)$. The tangent codimensions $(1,9,9,10)$ of
\cite[eq.~(8)]{Mo26}---the analogues of $(5,27,31,32)$ in his $25$-dimensional family of octics of bidegree
$(4,4)$, as read by machine---were reproduced modulo $p$ only. The octic program was validated on the quintic
example attributed to Maclean (two lines meeting in a point), which Theorem~A does not cover.

\section{Open points, and what a specialist should check first}\label{sec:open}

No human has read anything; which AI referee read which part is recorded in the Notice and in
Table~\ref{tab:status}. The parts that are new in version~2 were each read by exactly one AI referee, as working
notes; the proof of Lemma~\ref{lem:branch} was written by referee~8 and read by one later AI referee (referee~9)
only; version~2 as a whole was read cold by referee~9; version~3 has not been re-read.
\begin{enumerate}[label=\textup{(\arabic*)}]
\item \emph{The bridge at $r=-1$} (first point of referees~6, 7 and~8). The law $(r+1)/r$ rests on
$\theta_Q=y^2u'_Q$ (Lemma~\ref{lem:struct}) and on $B=A$ being read off from one family
(Lemmas~\ref{lem:Y}, \ref{lem:arc}, Proposition~\ref{prop:minus1}): the degenerate fibre $Y=Z_1\cup C'$, which is
not a complete intersection and is singular where $Z_1$ meets $\{h_1=0\}$; its cycle; flatness of the
two-parameter determinantal family; and (A) for it. It is the only non-formal geometric input, and
\eqref{eq:F} also uses it when (H) holds only for $Z_2$.
\item \emph{The standard facts in the generality used}: (A) for reducible and non-reduced fibres
(Lemma~\ref{lem:pi} for arbitrary $Z_i$; $Y$); the expansion \eqref{eq:exp} and (G1), i.e.\ that the natural
scheme structure of the locus of the non-primitive class $\delta_r$ is the zero scheme of the fixed-tube
integrals $\Phi_Q$, uniformly along $\VZ$; and that this is the analytic scheme considered in \cite{Mo26}.
\item \emph{Hypothesis \textup{(H)} and the scope of \eqref{eq:F}} (weakest points named by referee~7b).
(a) The vanishing \eqref{eq:I2} and the identity \eqref{eq:sss} are the whole content of Lemma~\ref{lem:R};
both rest on \eqref{eq:exp}, (A) and (G2). (b) (H) is proved only for lines, smooth conics, smooth plane curves
in the stated range and smooth $(2,2)$ curves; for all other types it is only computed, and the formula
$t_0=2\min+3\max-8$ is an observation modulo $p$: ``\eqref{eq:F} for every type'' is conditional. (c)
\eqref{eq:F} gives $A$ only on $(I(Z_1)I(Z_2))_{d-4}$, which is all of $K$ only when all $a_i,e_j\le3$; elsewhere
(e.g.\ $(1,1)+(e,e)$, $e\ge4$) the rest of $A$ is not described. Referee~7b did not check the statements on
Fermat surfaces of \S\ref{sec:comp}.
\item \emph{The $\Ob$-link} has no written proof (it uses the exact representative
$c\cdot\det\operatorname{Jac}$: period formula quoted from memory by referee~7, plus a flatness lemma). It is no
longer needed for $A\not\equiv0$, but the exact runs, the Lean certificate and all runs modulo $p$ depend on it.
No referee re-ran the exact programs or recompiled the Lean files.
\item \emph{Globalisation}: connectedness or irreducibility of $\cI^\circ$ and of $\mathcal U$; analyticity of
$\{A=0\}$; the bundle argument of Corollary~\ref{cor:L}(b); holomorphy of $E$ in the proof of Theorem~B(4); the
sign convention of $\cD$ (the order of the denominators in $\cD_{2j}$ was corrected by referee~6) and the
transcription of Lemma~\ref{lem:L} to the indices of Theorem~B, made for version~2.
\item \emph{Third order} is no longer needed for Theorem~B: part (4) follows from the linearity of $\cD$. The
direct third-order program at the Fermat point stays modulo $p$, one implementation, examined by no referee.
\item \emph{Reading of the literature.} \cite{Mo21}, \cite{Mo26} were read only by machine; a person should
check that family (9) of \cite{Mo21} and the slices of \cite{Mo26} are spanned by monomials (otherwise
Remark~\ref{rem:sym} does not explain the slice statements, and whether Theorem~B(5) is compatible with them
depends on the order of $u'$ along the slice, which we did not compute), and compare Conjecture~1, Proposition~4, the
abstract and Section~4 of \cite{Mo26}, and Theorem~1 and Remarks~6, 7 of \cite{Mo21}, with Table~\ref{tab:mov}
and \S\ref{sec:comp}. \cite[Thm~1.5]{Kl} (the index range of its inequality, and its clause on the excess; both
are now supported by a second machine reading, by referee~9, not by a person) and
\cite[\S5]{DFV} were read through machine digests, \cite{Kl2} by one agent; \cite{Mac}, \cite{Vo90},
\cite{Vo91} were not opened. The novelty statements rest on finite searches by the coordinating thread, referee~7
and a further agent.
\item \emph{Not covered}: $A$ outside $I(Z_1)I(Z_2)$ when some $a_i$ or $e_j$ is at least $4$; \textup{(H)},
hence $A\not\equiv0$, for all types; components of the Noether--Lefschetz locus of octics that do not meet
$\VZ$, and Harris's conjecture itself; $W_r$ beyond its germs at points of $\VZ$.
\item \emph{A common pattern, supported modulo $p$.} The hypothesis $d=a_1+a_2+e_1+e_2$ is the case of equality
in the inequality of \cite[Thm~1.5]{Kl} as we read it ($k=1$, $c=2$), and the configuration of \cite{P1} with
excess one (two linear subspaces with $(c-1)(d-2)=2$) is the case of equality for linear subspaces. In both, the
coefficient with unobstructed excess direction is $(-1)^{c-1}$. A modulo-$p$ program for two complete
intersections $Y_1=V(f_0,\dots,f_{c-1},p_c,\dots,p_k)$, $Y_2=V(h_0,\dots,h_{c-1},p_c,\dots,p_k)$ in a smooth
hypersurface of dimension $2k$ (file \path{gen_ci/gen_ci.py}; written by one AI agent, not refereed; two
cases re-run by the coordinating thread with a new seed) finds the same rule in every equality case tried that
has a nonzero obstruction space: for $k=2$, $d=4$, disjoint surfaces of types $(1,1,1)+(1,2,2)$,
$(1,1,2)+(1,1,2)$, $(1,1,1)+(1,1,3)$ give $\lambda_0=+1$ and surfaces sharing a quadric or a cubic equation give
$\lambda_0=-1$; seven types on quintic fourfolds and two on cubic sixfolds agree; in all, $35$ runs, $18$ types,
one to four random points each, with the validation cases included. Replacing an equation of $Y_i$ by its
partner in $F$ (a liaison) changes the sign of the period polynomial, changes $c$ by one and $\lambda$ to
$-\lambda$; this links several of these cases to linear ones and is consistent with the rule, but does not prove
it. Disjoint planes in a cubic fourfold are an equality case with no obstruction space, in which every $\lambda$
is unobstructed. A third manuscript \cite{P3} (AI-produced, refereed by AI agents only) proves the rule for
every $c$: with $F=f^{\mathsf T}Mh+\sum p_ls_l$ one has $[Y_1]+(-1)^{c-1}[Y_2]=m\,h^k-[R]$ on $X$, where
$R=V(p,M^{\mathsf T}f,\det M)$ is linked to $Y_1$ and lies on every member of the pencil $F-s\det M$; in the
range $(k-1)d<2k+2$ the second-order term is $\frac{\lambda-\lambda_0}{\lambda}A$. For $k=1$, $c=2$ this gives a
second description of $\NL(-1)$ in Theorem~A: near $X$ it is the set of surfaces containing a curve of the
form $V(M'^{\mathsf T}f',\det M')$, a curve of degree $(d-e_1)(d-e_2)-a_1a_2$ linked to a curve of the type of $Z_1$
in the complete intersection of the two forms $M'^{\mathsf T}f'$ (of degrees $d-e_1$, $d-e_2$). (This
item is the only change from version~3 to version~4; version~4 has not been re-read by a referee.)
\end{enumerate}

\begin{thebibliography}{99}\footnotesize\setlength{\itemsep}{0pt}\setlength{\parskip}{0pt}
\bibitem{BKU} G. Baldi, B. Klingler, E. Ullmo, \emph{Non-density of the exceptional components of the
Noether--Lefschetz locus}, arXiv:2312.11246; Int. Math. Res. Not. 2024, no.~21, 13642-- (search snippet only; data not verified).
\bibitem{CG} J. Carlson, P. Griffiths, \emph{Infinitesimal variations of Hodge structure and the global Torelli
problem}, 1980 (not opened; data unverified).
\bibitem{Dan1} A. Dan, \emph{Non-reduced components of the Noether--Lefschetz locus}, arXiv:1407.8491 (abstract
opened; partly machine-read).
\bibitem{Dan2} A. Dan, \emph{On a conjecture of Harris}, arXiv:2204.08079; Commun. Contemp. Math. (abstract
opened; machine-read).
\bibitem{DFV} J. Duque Franco, R. Villaflor Loyola, \emph{Periods of join algebraic cycles},
arXiv:2312.17222v4 (12 September 2025) (Section~5 read by the coordinating thread through machine digests of the
arXiv HTML).
\bibitem{Ei} D. Eisenbud, \emph{Commutative algebra with a view toward algebraic geometry}, Springer
(Hilbert--Burch; not opened).
\bibitem{Gr} P. Griffiths, \emph{On the periods of certain rational integrals I, II}, Ann. of Math. 90 (1969)
(not opened; data unverified).
\bibitem{Kl} R. Kloosterman, \emph{On a conjecture on Hodge loci of linear combinations of linear subvarieties},
arXiv:2312.12363v2; Rend. Circ. Mat. Palermo (2) 74 (2025), art.~187 (statements
machine-read; journal data as read by one agent on the publisher's page).
\bibitem{Kl2} R. Kloosterman, \emph{Hodge loci associated with linear subspaces intersecting in codimension
one}, arXiv:2401.10775; Math. Nachr., DOI 10.1002/mana.202400066 (machine-read by one agent; data not verified).
\bibitem{Mac} C. Maclean, \emph{A second-order invariant of the Noether--Lefschetz locus and two applications},
Asian J. Math. 9(3) (2005), 373--399 (not opened; data from \cite{Mo21}).
\bibitem{Mo21} H. Movasati, \emph{Special components of Noether--Lefschetz loci}, Rend. Circ. Mat. Palermo (2)
70 (2021), 861--874; arXiv:1908.04117 (abstract opened; body machine-read by one agent).
\bibitem{Mo26} H. Movasati, \emph{On a counterexample to a conjecture of J. Harris for octic surfaces},
arXiv:2606.26944 (25 June 2026) (abstract opened; body machine-read by three agents).
\bibitem{P1} \emph{Hodge loci of sums of two linear cycles: Movasati's conjecture for cubic sixfolds and
eightfolds}, AI-produced manuscript, version 5, 21 November 2026.
\bibitem{P3} \emph{The special coefficient for Hodge loci of sums of two complete intersections sharing
equations: a liaison identity and a second-order law}, AI-produced manuscript, 21 November 2026.
\bibitem{VL} R. Villaflor Loyola, \emph{Periods of complete intersection algebraic cycles}, Manuscripta Math. 167
(2022), 765--792; arXiv:1812.03964 (machine-read; used only for the computations).
\bibitem{Vo90} C. Voisin, \emph{Sur le lieu de Noether--Lefschetz en degr\'es 6 et 7}, Compositio Math. 75
(1990), 47--68 (not opened; data: \cite{Mo21}).
\bibitem{Vo91} C. Voisin, \emph{Contrexemple \`a une conjecture de J. Harris}, C. R. Acad. Sci. Paris 313 (1991),
685--687 (not opened; data: \cite{Mo21}).
\end{thebibliography}

\end{document}
